% concatenated from standard_gate_embeddings_v6.tex
\documentclass{article}%


%%%% Packages
\usepackage{graphicx}
\usepackage{multicol,multirow}
\usepackage{amsmath,amssymb,amsfonts}
\usepackage{mathrsfs}
\usepackage{amsthm}
\usepackage[margin=1in]{geometry}
\usepackage{rotating}
\usepackage{appendix}
\usepackage[numbers]{natbib}
\usepackage{ifpdf}
\usepackage[T1]{fontenc}
\usepackage{newtxtext}
\usepackage{newtxmath}
\usepackage{textcomp}
\usepackage{xcolor}
\definecolor{weblmcolor}{HTML}{1A5FB4}
\usepackage[colorlinks,allcolors=weblmcolor]{hyperref}
\usepackage{physics}

\newtheorem{theorem}{Theorem}[section]
\newtheorem{lemma}[theorem]{Lemma}
\newtheorem{proposition}[theorem]{Proposition}
\newtheorem{corollary}[theorem]{Corollary}
\theoremstyle{definition}
\newtheorem{definition}[theorem]{Definition}
\newtheorem{remark}[theorem]{Remark}
\newtheorem{example}[theorem]{Example}
\numberwithin{equation}{section}

\newcommand{\keywords}[1]{\par\noindent\textbf{Keywords:} #1}

%-----------------------
% Notation macros 
%-----------------------
\newcommand{\rA}{\mathsf{A}}
\newcommand{\rB}{\mathsf{B}}
\newcommand{\rI}{\mathsf{I}}
\newcommand{\rU}{\mathsf{U}}
\newcommand{\rV}{\mathsf{V}}
\newcommand{\rW}{\mathsf{W}}
\newcommand{\rH}{\mathsf{H}}
\newcommand{\rX}{\mathsf{X}}
\newcommand{\rY}{\mathsf{Y}}
\newcommand{\rZ}{\mathsf{Z}}
\newcommand{\R}{\mathbb{R}}
\newcommand{\N}{\mathbb{N}}
\newcommand{\Z}{\mathbb{Z}}
\newcommand{\C}{\mathbb{C}}
\newcommand{\Gr}{\mathrm{Gr}}
\newcommand{\End}{\mathrm{End}}
\newcommand{\Fix}{\mathrm{Fix}}
\newcommand{\Emb}{\mathrm{Emb}}
\newcommand{\Log}{\mathrm{Log}}
\newcommand{\diag}{\mathrm{diag}}
\newcommand{\Sym}{\mathrm{Sym}}
\newcommand{\Irr}{\mathrm{Irr}}
\newcommand{\Hom}{\mathrm{Hom}}
\newcommand{\St}{\mathrm{St}}
\newcommand{\Lie}{\mathrm{Lie}}
\newcommand{\Ran}{\mathrm{Ran}}

%\DeclareMathOperator{\tr}{tr}
\DeclareMathOperator{\Ad}{Ad}
%\DeclareMathOperator{\Tr}{Tr}


%\articletype{RESEARCH ARTICLE}
%\artid{20}
%\jyear{2021}
%\jvol{4}
%\jissue{1}
%\jdoi{10.1017/fms.2021.X}
%\raggedbottom

\title{Elementary Quantum Gates from Lie Group Embeddings in \texorpdfstring{$U(2^n)$}{U(2^n)}: Geometry, Universality, and Discretization}


\author{%
Antonio Falc\'o\\
Departamento de Matem\'aticas, F\'{\i}sica y Ciencias Tecnol\'ogicas\\
Universidad Cardenal Herrera-CEU, CEU Universities\\
San Bartolom\'e 55, 46115 Alfara del Patriarca (Valencia), Spain\\
\texttt{afalco@uchceu.es}
\and
Daniela Falc\'o--Pomares\\
Grupo de Investigaci\'on Bisite, Universidad de Salamanca\\
Calle Espejo s/n, 37007 Salamanca (Spain)\\
\texttt{dfp99@usal.es}
\and
Hermann G. Matthies\\
Institute of Scientific Computing\\
Technische Universit\"at Braunschweig, \\ 
Universit\"atsplatz 2, 38106 Braunschweig, Germany\\
\texttt{h.matthies@tu-bs.de}
}

\begin{document}

%\begin{Frontmatter}


%\date{}

%\received{04 May 2020}
%\revised{10 May 2020}
%\accepted{12 May 2020}

%\keywords[MSC Codes]{\codes[Primary]{81P68}; \codes[Secondary]{68Q12, 22E46}}

\maketitle

\begin{abstract}
In the standard circuit model, elementary gates are defined relative to a chosen tensor factorization and are therefore
extrinsic to the ambient group $U(2^n)$. Writing $N=2^n$, we introduce an \emph{intrinsic descriptor layer} in $U(N)$ by declaring as
primitive the motions inside faithful embedded copies of $SU(2)$ (phase-free), together with a phase-inclusive $U(2)$ variant.
We describe the embedding landscape $\Emb(SU(2),U(N))$ as a finite union of $U(N)$-homogeneous strata indexed by isotypic multiplicities,
with stabilizers given by centralizers, and we isolate a canonical \emph{two-level sector} parameterized by $\Gr_2(\C^N)$ up to a $PSU(2)$
gauge.
Equipping $U(N)$ with the Hilbert--Schmidt bi-invariant metric, each embedded subgroup is totally geodesic, yielding a variational
characterization of elementary motions via minimal-norm logarithms.
On the constructive side, we prove phase-free universality in $SU(N)$ from two-level primitives using QR/Givens factorizations together with
explicit diagonal generation, and we obtain full universality in $U(N)$ by explicit abelian phase bookkeeping (equivalently, via the $U(2)$
two-level dictionary).
Finally, we formalize a modular finite-alphabet compilation interface: any approximation routine in $SU(2)$ (e.g.\ Solovay--Kitaev) can be
lifted through two-level embeddings to yield $U(N)$-level synthesis with global operator-norm error control.
\end{abstract}

\keywords{quantum gate synthesis, quantum compiling, intrinsic gate descriptors, two-level unitaries (Givens/QR), $SU(2)$ embeddings in $U(2^n)$, Grassmannian, homogeneous manifolds, Riemannian geometry, universality, Solovay--Kitaev}



%\end{Frontmatter}

%\localtableofcontents

%\vspace*{14pt}
%============================================================
\section{Introduction}
\label{sec:introduction}
%============================================================

The circuit (gate) model of quantum computation is commonly presented after fixing a tensor
factorization $\mathbb{H}_{2^n}=(\C^2)^{\otimes n}$ and declaring as primitive those unitaries that
act on one or two tensor factors, from which larger circuits are obtained by composition
\cite{NielsenChuang2010}. Two foundational themes in this setting are \emph{expressiveness}
(universality) and \emph{synthesis} (compilation). Classical results show that mild hypotheses on
one- and two-qubit gate families yield dense generation of $U(2^n)$ up to phase
\cite{BarencoEtAl1995,DiVincenzo1995,DeutschBarencoEkert1995}, while constructive matrix-analytic
tooling---notably two-level (Givens/QR) factorizations---provides systematic routes to synthesis
\cite{ReckZeilingerBernsteinBertani1994}. For finite alphabets, the Solovay--Kitaev paradigm yields
polylogarithmic approximation lengths for dense generating sets
\cite{Kitaev2002,DawsonNielsen2006,HarrowRechtChuang2002}, and geometric viewpoints interpret circuit
complexity through geodesics and variational principles on unitary groups
\cite{NielsenDowlingGuDoherty2006,DowlingNielsen2008}.

Despite its operational success, the conventional notion of an ``elementary gate'' is \emph{extrinsic}
from the standpoint of the ambient group $U(2^n)$: it is tied to a chosen tensor layout and to
subsystem labels. Moreover, strict one-qubit locality alone does not generate the full computational
language for $n\ge 2$. This motivates an intrinsic question that is independent of hardware
addressing schemes and of basis choices:

\begin{quote}
\emph{Can one define elementarity directly inside $U(2^n)$ in a way that is conjugation-invariant and
does not presuppose a preferred tensor factorization?}
\end{quote}

\smallskip
\noindent\textbf{Embedding-based elementarity: an intrinsic descriptor layer.}
Fix $n\in\N$ and $N:=2^n$. We define elementarity intrinsically in $U(N)$ by declaring as primitive
the motions inside faithful \emph{embedded} copies of a two-level symmetry group. Concretely, we
model a logical two-level degree of freedom as a faithful embedding of $SU(2)$ into $U(N)$ and set
the phase-free dictionary
\[
\mathcal{G}^{SU}_{\mathrm{elem}}(n)
:=
\bigcup_{\phi\in \Emb(SU(2),U(N))}\phi\big(SU(2)\big)
\ \subset\ U(N).
\]
When one wishes to retain phases at the level of primitives, we also consider the $U(2)$ analogue
\[
\mathcal{G}^{U}_{\mathrm{elem}}(n)
:=
\bigcup_{\Phi\in \Emb(U(2),U(N))}\Phi\big(U(2)\big)
\ \subset\ U(N).
\]
This embedding-based notion is conjugation-invariant in $U(N)$, hence basis-independent, and it
separates the \emph{mathematical} definition of a primitive operation from any \emph{extrinsic}
addressing convention. Once a tensor structure is fixed, the usual single-qubit placements
$w^{(n)}_j:U(2)\hookrightarrow U(N)$ are recovered as particular embeddings, so tensor-local gates
embed into the intrinsic dictionaries; however tensor addressing is not built into the definition.

\smallskip
\noindent\textbf{The $SU(2)$/$U(2)$ split and phase management.}
The distinction between $SU(2)$ and $U(2)$ is structural and guides the entire paper. For fixed $N$,
$SU(2)$ has only finitely many unitary representation types in dimension $N$, which implies that the
conjugation action of $U(N)$ on $\Emb(SU(2),U(N))$ yields a \emph{finite} orbit-type stratification.
By contrast, $U(2)$ admits infinitely many determinant twists in the same dimension, so
$\Emb(U(2),U(N))$ naturally carries countably many orbit types even after fixing $N$. Accordingly,
our structural analysis is carried out at the $SU(2)$ level, while abelian diagonal/global phase
degrees of freedom are handled explicitly when passing from $SU(N)$ to $U(N)$.

\smallskip
\noindent\textbf{Main results and technical narrative.}
We develop a single pipeline from an intrinsic definition of elementarity in $U(N)$ to a modular synthesis/compilation interface.
At the structural level, we classify faithful $SU(2)$ embeddings in $U(N)$ up to conjugation, obtaining a finite orbit-type stratification
indexed by isotypic multiplicities and explicit stabilizers (centralizers).
At the computational level, we isolate the Grassmannian two-level sector (logical qubits as $W\in\Gr_2(\C^N)$ up to a $PSU(2)$ gauge),
prove constructive universality from two-level primitives via QR/Givens factorizations plus explicit diagonal generation, and formalize a
finite-alphabet interface by lifting $SU(2)$ approximation routines (e.g.\ Solovay--Kitaev) through two-level embeddings with global
operator-norm error control.

\smallskip
\noindent\textbf{Variational meaning.}
Equipping $U(N)$ with the Hilbert--Schmidt bi-invariant metric, we show that each embedded copy $\phi(G)\subset U(N)$
(with $G\in\{SU(2),U(2)\}$) is totally geodesic. Hence, within a fixed embedded gate manifold, energy-minimizing implementations of a target
element are exactly constant-speed one-parameter subgroups generated by minimal-norm logarithms in the embedded Lie algebra. In the two-level
sector this identifies the intrinsic metric cost in $U(N)$ with the corresponding $2\times 2$ generator, providing a direct variational
interpretation of elementary motions.

\subsection*{Positioning and novelty}
Several ingredients are classical in isolation: the representation theory of $SU(2)$, two-level (Givens/QR) synthesis, and finite-alphabet
approximation in $SU(2)$ (e.g.\ Solovay--Kitaev). The contribution here is to organize these components into a single \emph{intrinsic,
conjugation-invariant descriptor layer} for elementary gates internal to $U(2^n)$: primitives are motions inside faithful $SU(2)$/$U(2)$
embeddings, equipped with a natural gauge structure and parameter spaces (notably the Grassmannian two-level sector), and this viewpoint
yields a modular compilation interface obtained by lifting $SU(2)$ approximation routines through two-level embeddings with uniform
operator-norm control.



\medskip
\noindent\textbf{Contributions (intrinsic formalism and compilation interface).}
\smallskip
\begin{center}
\fbox{\parbox{\dimexpr\linewidth-2\fboxsep-2\fboxrule\relax}{
{\setlength{\leftmargini}{1.6em}\begin{itemize}\setlength{\itemsep}{2pt}
\item We introduce an \emph{intrinsic descriptor layer} for elementary operations in $U(N)$ ($N=2^n$) by declaring as primitives the motions inside faithful embedded copies of $SU(2)$ (and its $U(2)$ variant), yielding phase-free and phase-inclusive dictionaries.
\item We identify a canonical \emph{two-level sector} organized by $\Gr_2(\C^N)$ up to a $PSU(2)$ gauge, and we describe the embedding landscape as finitely many homogeneous strata indexed by isotypic multiplicities.
\item We prove \emph{phase-free universality} in $SU(N)$ from two-level primitives via QR/Givens factorization plus explicit diagonal generation; full universality in $U(N)$ follows by adjoining abelian diagonal/global factors (or by passing to the $U(2)$ two-level dictionary).
\item We provide a \emph{modular finite-alphabet compilation interface}: any $SU(2)$ approximation routine (e.g.\ Solovay--Kitaev) can be lifted through two-level embeddings with global operator-norm error control.
\end{itemize}}
}}
\end{center}

\noindent\textbf{Scope and non-claims.}
We do not claim improved worst-case complexity bounds for gate synthesis, nor optimality of any particular compilation subroutine.
Our aim is to supply an intrinsic, representation- and gauge-aware language for elementary gates in $U(2^n)$ together with a clean, modular
pipeline connecting $U(N)$-factorizations to $SU(2)$ finite-alphabet approximation.

\medskip
\noindent\textbf{Supporting layers (scope alignment).}
Sections~\ref{sec:EmbSU2UN} and~\ref{sec:variational_principle} provide structural and interpretive backbones for the intrinsic descriptor
layer, rather than additional compilation claims. The embedding-landscape analysis makes the admissible intrinsic two-level degrees of
freedom \emph{computable} (finite orbit types with explicit stabilizers/centralizers), while the variational section provides a clean
baseline metric cost for elementary motions (total geodesy and minimal-logarithm characterizations) that can be replaced by more
hardware-specific control models.


\paragraph{Organization of the paper.}
Section~\ref{sec:preliminaries} fixes notation and formalizes the embedding-based dictionaries,
including the $SU(2)$/$U(2)$ phase split and an intrinsic locality baseline used for comparison.
Section~\ref{sec:logical_qubits} isolates the canonical two-level stratum, identifies it with the
Grassmannian $\Gr_2(\C^N)$, makes the associated $PSU(2)$ gauge precise, and relates this intrinsic
model to conventional tensor-factor locality. Section~\ref{sec:universality} proves phase-free
universality at the two-level level and explains how full universality in $U(N)$ is obtained by
explicit diagonal/global-phase management. Section~\ref{sec:SK} develops the finite-alphabet
interface by combining two-level factorizations with the Solovay--Kitaev paradigm on $SU(2)$ and
lifting the resulting words to $U(N)$ with global operator-norm error control.
Section~\ref{sec:EmbSU2UN} provides the structural backbone by describing the conjugation action of
$U(N)$ on $\Emb(SU(2),U(N))$, identifying stabilizers via centralizers, and deriving the finite
orbit-type decomposition indexed by irreducible multiplicity data, together with explicit dimension
formulas and representative examples. Section~\ref{sec:variational_principle} establishes the
variational characterization of elementary motions via total geodesy of embedded subgroups.
Finally, Section~\ref{sec:conclusion} summarizes the main results and indicates natural directions
for further work.

%============================================================
\section{Preliminary definitions and results}
\label{sec:preliminaries}
%============================================================

\noindent
This section fixes notation and conventions used throughout, and collects the auxiliary facts needed
to read the main results with minimal cross-referencing. Our viewpoint is intrinsic: we work inside
the ambient unitary group $U(N)$, $N=2^n$, and define elementary operations as motions inside faithful
embedded copies of a two-level symmetry group.

\smallskip
\noindent\textbf{Reader's guide.}
For a first pass through the main narrative (canonical two-level sector $\rightarrow$ universality
$\rightarrow$ finite-alphabet compilation), it suffices to skim
Section~\ref{subsec:hilbert_and_groups} and keep in mind the embedding-based dictionaries introduced in
Section~\ref{subsec:elementary_via_embeddings}. The Hilbert--Schmidt geometry and generator
normalizations fixed in Sections~\ref{subsec:lie_algebra_hs}--\ref{subsec:pauli_strings} are used
primarily in the variational interpretation of Section~\ref{sec:variational_principle}. Finally,
Section~\ref{subsec:intrinsic_locality} provides an intrinsic commutant-based notion of tensor
locality used only as a baseline comparison with our embedding-elementary framework.

\smallskip
\noindent\textbf{The $SU(2)$/$U(2)$ split.}
We take $U(N)$ as ambient group so that global phases and diagonal factors remain explicit in
synthesis and compilation statements. At the same time, the non-abelian two-level motions that drive
our orbit stratification and discretization are naturally governed by embeddings of $SU(2)$ into
$U(N)$. Accordingly, the structural analysis is carried out at the $SU(2)$ level (yielding finitely
many orbit types for fixed $N$), while abelian $U(1)$ phases and diagonal factors are managed
separately, either by canonical normalization to $SU(N)$ or via an explicit diagonal-compilation
step.


%------------------------------------------------------------
\subsection{Multi-qubit Hilbert space, ambient groups, and global-phase conventions}
\label{subsec:hilbert_and_groups}
Fix $n\in\N$ and set $N:=2^n$. We consider the $n$-qubit Hilbert space
\begin{equation}\label{eq:HN_def}
\mathbb{H}_N := (\mathbb{C}^2)^{\otimes n} \cong \mathbb{C}^{N},
\end{equation}
equipped with the standard Hermitian inner product. Once a computational basis is fixed, we identify
linear operators on $\mathbb{H}_N$ with matrices in $\mathbb{C}^{N\times N}$.

\paragraph{Global phases and the choice of ambient group.}
In many quantum-information settings one works modulo global phase, i.e.\ in the projective unitary
group $PU(N)=U(N)/U(1)$ (or, equivalently, within $SU(N)$). In the present work we keep global phases
explicit and take $U(N)$ as ambient Lie group. This is deliberate: our central object is the space
of embeddings $U(2)\hookrightarrow U(N)$, which is naturally defined at the level of $U(N)$ without
imposing determinant-one constraints, and retaining phases allows us to state synthesis and
compilation procedures uniformly.

When a phase-free formulation is convenient, we pass canonically to $SU(N)$ by removing the
determinant phase:
\begin{equation}\label{eq:UtoSU}
\widetilde{U} := e^{-i\arg(\det U)/N}\,U \in SU(N).
\end{equation}
Note that $U$ and $\widetilde{U}$ induce the same element of $PU(N)$, so this operation changes only
the global phase and not the projective action on pure states.


%------------------------------------------------------------
\subsection{Lie algebra, Hilbert--Schmidt geometry, and the exponential map}
\label{subsec:lie_algebra_hs}
Let $\mathfrak{u}(N)$ denote the Lie algebra of $U(N)$, i.e.\ the space of anti-Hermitian matrices:
\[
\mathfrak{u}(N)=\{X\in\mathbb{C}^{N\times N}: X^\dagger=-X\}.
\]
We equip $\mathfrak{u}(N)$ with the Hilbert--Schmidt pairing
\begin{equation}\label{eq:HS_pairing}
\langle X,Y\rangle_{\mathrm{hs}} := \frac12 \Tr(X^\dagger Y),\qquad X,Y\in\mathfrak{u}(N).
\end{equation}
This pairing is $\Ad$-invariant and induces the standard bi-invariant Riemannian metric on $U(N)$.
For a smooth curve $U:[0,1]\to U(N)$, the \emph{body velocity} (right-trivialized velocity) is
\begin{equation}\label{eq:body_velocity}
A(t):=\dot U(t)\,U(t)^{-1}\in\mathfrak{u}(N),
\end{equation}
and the corresponding energy functional is
\begin{equation}\label{eq:energy_hs}
\mathcal{E}[U]:=\frac12\int_0^1 \|A(t)\|_{\mathrm{hs}}^2\,dt
=\frac12\int_0^1 \|\dot U(t)\|^2\,dt.
\end{equation}
With respect to this bi-invariant metric, geodesics in $U(N)$ are precisely one-parameter subgroups
\begin{equation}\label{eq:geodesics_one_parameter}
U(t)=U(0)\exp(tX),\qquad X\in\mathfrak{u}(N).
\end{equation}
We adopt the matrix exponential $\exp:\mathfrak{u}(N)\to U(N)$ as the default exponential map.


%------------------------------------------------------------
\subsection{Pauli matrices, Pauli strings, and normalization of generators}
\label{subsec:pauli_strings}
Let $\{\sigma_0,\sigma_x,\sigma_y,\sigma_z\}$ denote the Pauli matrices, with $\sigma_0:=I_2$.
A \emph{Pauli string} is an operator of the form
\begin{equation}\label{eq:pauli_string}
P=\sigma_{a_1}\otimes\cdots\otimes\sigma_{a_n},\qquad a_j\in\{0,x,y,z\}.
\end{equation}
Pauli strings form an orthogonal basis of $\mathbb{C}^{N\times N}$ under the Hilbert--Schmidt inner
product and provide a convenient coordinate system for describing embedded $\mathfrak{su}(2)$
directions.

We associate to each nontrivial Pauli string $P\neq I_N$ the normalized anti-Hermitian generator
\begin{equation}\label{eq:PauliGenerator}
T_P := -\frac{i}{2}\,P \in \mathfrak{u}(N).
\end{equation}
This convention matches the standard parametrization of single-qubit rotations
$R_\alpha(\theta)=\exp(-i\theta\sigma_\alpha/2)$ and fixes the sign and the factor $1/2$ appearing in
all elementary exponentials. In particular, if $P=\sigma_\alpha^{(j)}$ is a single-site Pauli
operator acting on the $j$th qubit, then $\exp(\theta T_P)$ is a rotation of angle $\theta$ around
axis $\alpha$ on the corresponding degree of freedom.


%------------------------------------------------------------
\subsection{Intrinsic locality via commutants and the canonical reference subclass}
\label{subsec:intrinsic_locality}
To benchmark embedding-based elementarity against conventional tensor locality, we use an intrinsic
(commutant-based) definition of ``acting on a subset of qubits'' relative to the fixed
factorization~\eqref{eq:HN_def}. Fix $S\subset\{1,\dots,n\}$ and write $\bar S$ for its complement.
Using the tensor-ordering convention fixed throughout (equivalently, a canonical permutation unitary
$\Pi_S$ that brings the tensor factors in $S$ to the front), define the subgroup of unitaries acting
trivially on $S$ by
\begin{equation}\label{eq:Ubar_def}
\mathcal{U}_{\bar S}:=\Pi_S^{-1}\big(I_{2^{|S|}}\otimes U(2^{n-|S|})\big)\Pi_S \subset U(N).
\end{equation}

\begin{definition}[$S$-local unitaries via commutants]\label{def:S_local_commutant}
The \emph{$S$-local unitary subgroup} is the commutant of $\mathcal{U}_{\bar S}$ in $U(N)$,
\begin{equation}\label{eq:S_local_commutant}
U(N)_S:=\{\,U\in U(N): UX=XU\ \ \forall X\in \mathcal{U}_{\bar S}\,\}.
\end{equation}
\end{definition}

\noindent
By the bicommutant theorem (or direct matrix-algebra computations), one has the canonical
identification
\begin{equation}\label{eq:S_local_structure}
U(N)_S = \Pi_S^{-1}\big(U(2^{|S|})\otimes I_{2^{n-|S|}}\big)\Pi_S.
\end{equation}
Thus $U(N)_S$ recovers the usual tensor-factor notion of acting nontrivially only on the degrees of
freedom in $S$, but expressed intrinsically through commutation relations.

\paragraph{Canonical one-qubit reference subclass.}
For $S=\{j\}$ one obtains the one-site local subgroup $U(N)_{\{j\}}$. We define the strictly local
one-qubit family
\begin{equation}\label{eq:Gloc_def}
\mathcal{G}_{\mathrm{loc}}(n):=\bigcup_{j=1}^n U(N)_{\{j\}},
\end{equation}
which serves as the canonical reference subclass representing the conventional one-qubit gate
paradigm. In particular, $\mathcal{G}_{\mathrm{loc}}(n)$ will be used as a baseline when comparing
tensor-local elementarity to the embedding-elementary dictionaries defined next.


%------------------------------------------------------------
\subsection{Elementary gates via Lie group embeddings}
\label{subsec:elementary_via_embeddings}
The commutant-based notion of locality depends on the fixed tensor factorization~\eqref{eq:HN_def}.
Our proposed elementary-gate framework adopts a representation-theoretic viewpoint: a two-dimensional
logical degree of freedom is treated as primitive, independently of a prescribed tensor layout, by
considering faithful realizations of $U(2)$ (or, for stratification purposes, $SU(2)$) inside $U(N)$.

\medskip
\noindent\textbf{Embeddings and gate manifolds.}
Let $G\in\{SU(2),U(2)\}$. We write $\Emb(G,U(N))$ for the set of faithful Lie group embeddings
$\phi:G\hookrightarrow U(N)$. Each $\phi$ identifies a concrete two-level gate manifold
$\phi(G)\subset U(N)$ and an embedded Lie algebra $d\phi_e(\mathfrak g)\subset\mathfrak u(N)$.

\begin{definition}[Lie group embeddings]\label{def:embedding}
Let $N\in\mathbb{N}$ and let $G\in\{SU(2),U(2)\}$.
We define the set of Lie group embeddings of $G$ into $U(N)$ by
\[
\Emb(G,U(N))
:=
\left\{\phi:G\to U(N)\ \middle|\
\begin{array}{l}
\phi \text{ is a continuous group homomorphism and injective,}\\[2pt]
d\phi_e:\mathfrak g\to\mathfrak{u}(N)\text{ is injective}
\end{array}
\right\},
\]
where $\mathfrak g$ denotes the Lie algebra of $G$.
Since $G$ is compact, any continuous homomorphism $G\to U(N)$ is automatically smooth. Moreover,
injectivity implies that $\phi$ is a topological embedding whose image $\phi(G)$ is a closed Lie
subgroup of $U(N)$, and $d\phi_e$ identifies $\mathfrak g$ with the Lie algebra of $\phi(G)$.
\end{definition}

\medskip
\noindent\textbf{Embedding-elementary dictionaries.}
In view of Remark~\ref{rem:U2_vs_SU2}, it is useful to distinguish the phase-free non-abelian content
from the full unitary group.

\paragraph{Convention (two elementary dictionaries).}
Let $n\in\mathbb{N}$ and $N:=2^n$. We define
\begin{align}
\label{eq:Gelem_def_SU}
\mathcal{G}^{SU}_{\mathrm{elem}}(n)
&:=\bigcup_{\phi\in \Emb(SU(2),U(N))}\phi\big(SU(2)\big)\ \subset\ U(N),
\\
\label{eq:Gelem_def_U}
\mathcal{G}^{U}_{\mathrm{elem}}(n)
&:=\bigcup_{\phi\in \Emb(U(2),U(N))}\phi\big(U(2)\big)\ \subset\ U(N).
\end{align}
By construction, both $\mathcal{G}^{SU}_{\mathrm{elem}}(n)$ and $\mathcal{G}^{U}_{\mathrm{elem}}(n)$
are invariant under conjugation by $U(N)$.

The set $\mathcal{G}^{SU}_{\mathrm{elem}}(n)$ encodes the non-abelian elementary motions (with global
phases and diagonal factors treated separately). In Section~\ref{sec:universality} we show that the
subgroup it generates contains (indeed generates) $SU(N)$ via two-level factorizations. When full
$U(N)$ reachability is required, we work instead with $\mathcal{G}^{U}_{\mathrm{elem}}(n)$, or
equivalently supplement the $SU(2)$-based synthesis with explicit compilation of the abelian factors
(global phase and trailing diagonal terms) as in Section~\ref{sec:SK}.

\paragraph{Canonical local layer (reference subclass).}
The embedding-based framework is intrinsically defined on $U(N)$ and does not require a preferred
tensor factorization of $\mathbb{H}_N$. Nevertheless, once the identification
$\mathbb{H}_N\cong(\mathbb{C}^2)^{\otimes n}$ is fixed, there are canonical ``coordinate'' embeddings
\[
w_j^{(n)}:U(2)\hookrightarrow U(N),\qquad
w_j^{(n)}(V)=I_2^{\otimes (j-1)}\otimes V\otimes I_2^{\otimes (n-j)},\qquad j=1,\dots,n,
\]
whose images coincide with the commutant-defined one-site subgroups $U(N)_{\{j\}}$ in
\eqref{eq:S_local_structure}. Equivalently,
\[
\mathcal{G}_{\mathrm{loc}}(n):=\bigcup_{j=1}^n w_j^{(n)}\big(U(2)\big)\subset U(N),
\qquad
\bigcup_{j=1}^n w_j^{(n)}\big(SU(2)\big)\subset SU(N)
\]
are the canonical local layers singled out by the chosen tensor structure.

\begin{remark}[Why we stratify by $SU(2)$-embeddings]\label{rem:U2_vs_SU2}
Restricting $\phi\in\Emb(U(2),U(N))$ to $SU(2)\subset U(2)$ yields an element of $\Emb(SU(2),U(N))$.
Conversely, an embedding of $SU(2)$ captures the full non-abelian content of a logical two-level
degree of freedom. For fixed $N$, unitary representation types of $U(2)$ are infinite (determinant
twists), whereas those of $SU(2)$ are finite. Since our discretization step synthesizes $SU(2)$
motions (with global phases and diagonal factors treated separately), we formulate the orbit
stratification and the intrinsic ``logical-qubit'' degrees of freedom using $\Emb(SU(2),U(N))$.
\end{remark}


%------------------------------------------------------------
\subsection{Embeddings into \texorpdfstring{$U(N)$}{U(N)} and faithful unitary representations}
\label{subsec:embeddings_representations}
The embedding-based notion of elementarity admits a convenient reformulation in finite-dimensional
representation theory. Any continuous homomorphism $\phi:G\to U(N)$ prescribes an action of $G$ on
$\C^N$ by unitary operators and can therefore be viewed as an $N$-dimensional unitary representation.
The embedding requirement (injectivity and identification with a Lie subgroup) amounts to
\emph{faithfulness} (trivial kernel), while conjugation in $U(N)$ corresponds to unitary change of
basis on $\C^N$. This perspective supplies discrete invariants (isotypic multiplicities), clarifies
the orbit structure of the $U(N)$-action on $\Emb(G,U(N))$, and underlies the homogeneous
stratification developed in Section~\ref{sec:EmbSU2UN}; see, e.g., \cite{Hall2015,BrockerDieck1985,Knapp2002}.

\begin{definition}[Unitary representation]\label{def:unitary_representation}
A \emph{unitary representation} of a topological group $G$ on $\C^N$ is a continuous homomorphism
$\rho:G\to U(N)$. It is \emph{faithful} if $\ker(\rho)=\{e\}$.
\end{definition}

\noindent
With Definition~\ref{def:unitary_representation}, any continuous homomorphism $\phi:G\to U(N)$ (with
$G$ compact) is a finite-dimensional unitary representation; moreover, injectivity of $\phi$ is
equivalent to faithfulness of the associated representation. We will apply this both to $G=U(2)$
(when tracking central phases) and to $G=SU(2)$ (for the orbit stratification).

\begin{proposition}[Embeddings $\boldsymbol{G\hookrightarrow U(N)}$ as faithful unitary representations]
\label{prop:embeddings_as_faithful_representations_G}
Fix $N\in\mathbb{N}$ and let $G\in\{U(2),SU(2)\}$. There is a natural identification between Lie group
embeddings $\phi\in \Emb(G,U(N))$ and faithful $N$-dimensional unitary representations of $G$. More precisely:
\begin{enumerate}
\item[(i)] Any continuous homomorphism $\phi:G\to U(N)$ is a unitary representation on $\C^N$; we denote it by $\rho_\phi:=\phi$.
Moreover, $\phi$ is injective if and only if $\rho_\phi$ is faithful.
\item[(ii)] Conversely, any faithful continuous unitary representation $\rho:G\to U(N)$ defines an element of $\Emb(G,U(N))$.
\item[(iii)] The conjugation action of $U(N)$ on $\Emb(G,U(N))$,
\[
(\Ad_W\cdot \phi)(g):=W\,\phi(g)\,W^{-1}\qquad (W\in U(N)),
\]
corresponds exactly to unitary equivalence of representations.
\end{enumerate}
\end{proposition}

\begin{proof}
(i) By Definition~\ref{def:unitary_representation}, any continuous homomorphism $\phi:G\to U(N)$ is a unitary representation on $\C^N$,
which we denote by $\rho_\phi:=\phi$. Since $\ker(\rho_\phi)=\ker(\phi)$, the representation $\rho_\phi$ is faithful if and only if
$\phi$ is injective.

(ii) Let $\rho:G\to U(N)$ be a faithful continuous unitary representation. Since $G$ is compact and $U(N)$ is Hausdorff, $\rho(G)$ is
compact and therefore closed in $U(N)$. Hence $\rho(G)\le U(N)$ is a closed subgroup and, by the closed subgroup theorem, an embedded
Lie subgroup. Because $\rho$ is continuous and injective, $\rho:G\to\rho(G)$ is a homeomorphism. Finally, since $G$ is a compact Lie
group, any continuous homomorphism $G\to U(N)$ is automatically smooth (see, e.g., \cite[Ch.~4]{Hall2015}); hence $\rho$ is a Lie
group embedding, i.e.\ $\rho\in\Emb(G,U(N))$.

(iii) If $\phi'=\Ad_W\cdot\phi$ for some $W\in U(N)$, then $\rho_{\phi'}(g)=W\rho_\phi(g)W^{-1}$ for all $g\in G$, so $\rho_{\phi'}$
and $\rho_\phi$ are unitarily equivalent. Conversely, if two unitary representations $\rho,\rho':G\to U(N)$ are unitarily equivalent,
then there exists $W\in U(N)$ such that $\rho'(g)=W\rho(g)W^{-1}$ for all $g\in G$, i.e.\ $\rho'=\Ad_W\cdot\rho$.
\end{proof}

\noindent
Proposition~\ref{prop:embeddings_as_faithful_representations_G} reduces the study of $\Emb(G,U(N))$
(for $G\in\{U(2),SU(2)\}$) to the classification of finite-dimensional unitary representations of the
compact group $G$. In particular, unitary representations of compact groups are completely reducible,
and their decomposition into irreducibles is unique up to unitary equivalence; hence conjugacy
classes of embeddings are encoded by discrete multiplicity data in the isotypic decomposition (see,
e.g., \cite{BrockerDieck1985,Hall2015}).

The next corollary makes this statement explicit and will serve as representation-theoretic input
for the homogeneous decomposition results in Section~\ref{sec:EmbSU2UN}.

\begin{corollary}[Conjugacy classes of embeddings and isotypic multiplicity data]
\label{cor:conjugacy_isotypic_data_G}
Fix $N\in\mathbb{N}$ and let $G\in\{U(2),SU(2)\}$. Two embeddings $\phi_1,\phi_2\in \Emb(G,U(N))$ are
conjugate in $U(N)$ if and only if the associated faithful unitary representations $\rho_{\phi_1},
\rho_{\phi_2}$ are unitarily equivalent. Equivalently, $\rho_{\phi_1}$ and $\rho_{\phi_2}$ have the
same isotypic decomposition (the same irreducible multiplicity data).

More precisely, there exist finite-support multiplicity families $(m_\lambda^{(1)})_\lambda$ and
$(m_\lambda^{(2)})_\lambda$ and $G$-equivariant unitary isomorphisms
\[
\C^N \ \cong\ \bigoplus_{\lambda}\Big(\C^{m_\lambda^{(k)}}\otimes V_\lambda\Big),
\qquad k\in\{1,2\},
\]
where $\{V_\lambda\}_\lambda$ ranges over a complete set of representatives of the unitary
equivalence classes of finite-dimensional irreducible unitary representations of $G$. In
particular, $\phi_1$ and $\phi_2$ are conjugate in $U(N)$ if and only if
\[
m_\lambda^{(1)}=m_\lambda^{(2)} \quad \text{for all }\lambda.
\]
\end{corollary}

\begin{proof}
By Proposition~\ref{prop:embeddings_as_faithful_representations_G}(iii), $\phi_1$ and $\phi_2$ are
conjugate in $U(N)$ if and only if $\rho_{\phi_1}$ and $\rho_{\phi_2}$ are unitarily equivalent.

Since $G$ is compact, every finite-dimensional unitary representation of $G$ is completely reducible:
it decomposes as a finite orthogonal direct sum of irreducibles (unitary Maschke theorem); see,
e.g., \cite{BrockerDieck1985,Hall2015,Knapp2002}. Moreover, the unitary equivalence class of a
finite-dimensional unitary representation is uniquely determined by its irreducible multiplicities,
e.g.\ via character orthogonality (equivalently, Schur orthogonality for matrix coefficients); see,
e.g., \cite{BrockerDieck1985,Hall2015}. Therefore $\rho_{\phi_1}\simeq \rho_{\phi_2}$ if and only if
$m_\lambda^{(1)}=m_\lambda^{(2)}$ for all $\lambda$, and this criterion is equivalent to conjugacy of
$\phi_1,\phi_2$ in $U(N)$.
\end{proof}


%============================================================
\section{Logical Qubits and the Grassmannian Model}
\label{sec:logical_qubits}
%============================================================

\noindent
The embedding-based framework treats a two-dimensional logical degree of freedom as primitive,
independently of any preferred tensor factorization of $\mathbb{H}_N\cong\C^N$. In this section we
isolate a distinguished and computationally transparent family of embeddings---the \emph{two-level
embeddings}---and show that they are naturally organized by the complex Grassmannian $\Gr_2(\C^N)$.

\smallskip
\noindent
\textbf{Takeaway.} A \emph{two-level logical qubit} is a two-plane $W\in\Gr_2(\C^N)$ (support data),
and a two-level gate is an $SU(2)$ action on $W$ extended by the identity on $W^\perp$. Choosing a
frame on $W$ produces an explicit embedding $SU(2)\hookrightarrow U(N)$, but the resulting embedding
data is defined only up to the inner-automorphism gauge of $SU(2)$, i.e.\ up to $PSU(2)\cong SO(3)$.
This Grassmannian sector will serve as the canonical computational stratum of the full embedding
landscape.

%------------------------------------------------------------
\subsection{Two-level subgroups and two-level embeddings}
\label{subsec:two_level_embeddings}
%------------------------------------------------------------

\begin{definition}[Two-level unitary and special-unitary subgroups]\label{def:two_level_subgroup}
Let $W\subset \C^N$ be a complex subspace with $\dim_\C W=2$, and let $W^\perp$ be its orthogonal
complement. Define the \emph{two-level unitary subgroup supported on $W$} by
\begin{equation}\label{eq:two_level_subgroup_UN}
U(N)[W]\;:=\;\{\,U\in U(N): U|_{W}\in U(W)\ \text{and}\ U|_{W^\perp}=I_{W^\perp}\,\},
\end{equation}
and the \emph{two-level special-unitary subgroup supported on $W$} by
\begin{equation}\label{eq:two_level_subgroup_SUN}
SU(N)[W]\;:=\;\{\,U\in U(N): U|_{W}\in SU(W)\ \text{and}\ U|_{W^\perp}=I_{W^\perp}\,\}.
\end{equation}
Equivalently, $U(N)[W]$ acts arbitrarily on $W$ and trivially on $W^\perp$, while $SU(N)[W]$ imposes
$\det(U|_W)=1$.
\end{definition}

\begin{remark}[Non-canonical identifications]
The groups $U(N)[W]$ and $SU(N)[W]$ are (non-canonically) isomorphic to $U(2)$ and $SU(2)$,
respectively. Making such an identification amounts to choosing a unitary isomorphism $W\cong\C^2$,
i.e.\ an orthonormal frame of $W$.
\end{remark}

\begin{definition}[Framed two-level $SU(2)$-embedding]\label{def:framed_two_level_embedding}
Let $W\in\Gr_2(\C^N)$ and let $f:\C^2\to W$ be a unitary isomorphism (an orthonormal frame on $W$).
Define
\begin{equation}\label{eq:two_level_embedding_SU2}
\phi_{W,f}:SU(2)\hookrightarrow U(N),\qquad
\phi_{W,f}(S)\;:=\; f S f^{-1}\ \oplus\ I_{W^\perp},
\end{equation}
where the direct sum is taken with respect to $\C^N=W\oplus W^\perp$.
\end{definition}

\begin{proposition}[Image and gauge covariance]\label{prop:gauge_invariance_SU2}
For fixed $W\in\Gr_2(\C^N)$, the subgroup $\phi_{W,f}(SU(2))\subset U(N)$ is independent of the
choice of frame $f$ and equals $SU(N)[W]$. More precisely, if $f,f':\C^2\to W$ are unitary
isomorphisms, then there exists $S_0\in SU(2)$ such that
\[
\phi_{W,f'}(S)=\phi_{W,f}(S_0 S S_0^{-1}),\qquad S\in SU(2),
\]
hence $\phi_{W,f'}(SU(2))=\phi_{W,f}(SU(2))=SU(N)[W]$.
\end{proposition}

\begin{proof}
Since $f,f'$ are unitary isomorphisms onto $W$, the map $U:=f^{-1}f'$ lies in $U(2)$ and $f'=fU$.
Let $\delta:=\det(U)\in U(1)$ and set $S_0:=\delta^{-1/2}U\in SU(2)$ (choose any square root
$\delta^{1/2}$; the ambiguity is $\pm I$ and does not affect the inner action). Then $USU^{-1}=S_0SS_0^{-1}$
for all $S\in SU(2)$, because the central scalar cancels. Therefore
\[
\phi_{W,f'}(S)=(fU)S(U^{-1}f^{-1})\oplus I_{W^\perp}
=f(USU^{-1})f^{-1}\oplus I_{W^\perp}
=\phi_{W,f}(S_0SS_0^{-1}).
\]
The equality of images follows, and $\phi_{W,f}(SU(2))$ acts as $SU(W)$ on $W$ and fixes $W^\perp$
pointwise, hence equals $SU(N)[W]$.
\end{proof}


%------------------------------------------------------------
\subsection{The Grassmannian as the configuration space of two-level logical qubits}
\label{subsec:grassmannian_logical_qubits}
%------------------------------------------------------------

The construction above associates to each two-plane $W$ a canonical two-level subgroup $SU(N)[W]$.
Conversely, a subgroup of the form $SU(N)[W]$ uniquely determines $W$ as the maximal subspace on which
it can act nontrivially.

\begin{proposition}[Two-level $SU(2)$ subgroups are parametrized by $\Gr_2(\C^N)$]
\label{prop:two_level_param_by_grass}
The assignment
\begin{equation}\label{eq:W_to_subgroup}
\Psi:\Gr_2(\C^N)\longrightarrow \{\text{two-level $SU(2)$ subgroups of }U(N)\},\qquad
\Psi(W):=SU(N)[W],
\end{equation}
is injective. In particular, two-level embedded copies of $SU(2)$ in $U(N)$ are naturally
parametrized by $\Gr_2(\C^N)$.
\end{proposition}

\begin{proof}
If $SU(N)[W_1]=SU(N)[W_2]$, then their fixed-point subspaces coincide:
\[
\Fix\big(SU(N)[W_i]\big)\;:=\;\bigcap_{U\in SU(N)[W_i]}\ker(U-I)=W_i^\perp.
\]
Indeed, every element fixes $W_i^\perp$ pointwise, while the restriction to $W_i$ ranges over
$SU(W_i)$ and hence has no nonzero fixed vector. Therefore $W_1^\perp=W_2^\perp$ and $W_1=W_2$.
\end{proof}

\begin{definition}[Two-level logical qubit]\label{def:logical_qubit}
A \emph{(two-level) logical qubit} in $\C^N$ is a choice of a complex $2$-plane $W\in\Gr_2(\C^N)$.
A \emph{framed logical qubit} is a pair $(W,f)$ where $W\in\Gr_2(\C^N)$ and $f:\C^2\to W$ is a unitary
isomorphism.
\end{definition}

\begin{remark}[Gauge group and Stiefel-bundle viewpoint]\label{rem:gauge_bundle_viewpoint}
Framed logical qubits form the Stiefel manifold $\St_2(\C^N)$ of orthonormal $2$-frames, and the
projection $\St_2(\C^N)\to \Gr_2(\C^N)$ is a principal $U(2)$-bundle. However, for the associated
$SU(2)$-embedding \eqref{eq:two_level_embedding_SU2}, the central $U(1)\subset U(2)$ acts trivially:
if $f'=fe^{i\theta}I$ then $\phi_{W,f'}=\phi_{W,f}$. Accordingly, the effective gauge acting on the
embedding data is
\[
U(2)/U(1)\ \cong\ SU(2)/\{\pm I\}\ =:\ PSU(2)\ \cong\ SO(3),
\]
i.e.\ the group of inner automorphisms of $SU(2)$.
\end{remark}


%------------------------------------------------------------
\subsection{The two-level stratum as a homogeneous space and a principal bundle}
\label{subsec:two_level_stratum_bundle}
%------------------------------------------------------------

The family of framed embeddings $(W,f)\mapsto \phi_{W,f}$ defines a distinguished subset of
$\Emb(SU(2),U(N))$, corresponding to the representation type
$V_{1}\oplus V_{0}^{\oplus(N-2)}$ (cf.\ Remark~\ref{rem:two_level_stratum_concise} and the general
classification in Section~\ref{sec:EmbSU2UN}). This subset is a single $U(N)$-conjugacy class, and
its internal redundancy is precisely the $PSU(2)$ gauge discussed above.

\begin{proposition}[Two-level embeddings form a principal $PSU(2)$-bundle over $\Gr_2(\C^N)$]
\label{prop:two_level_principal_bundle_SU2}
Let $\Emb_{\mathrm{2lvl}}(SU(2),U(N))\subset \Emb(SU(2),U(N))$ denote the set of embeddings of the form
$\phi_{W,f}$ in \eqref{eq:two_level_embedding_SU2}. Then:
\begin{enumerate}
\item[(i)] The conjugation action of $U(N)$ is transitive on $\Emb_{\mathrm{2lvl}}(SU(2),U(N))$; hence
this set is a smooth homogeneous manifold (the \emph{two-level stratum}).
\item[(ii)] The map
\[
\pi:\Emb_{\mathrm{2lvl}}(SU(2),U(N))\to \Gr_2(\C^N),\qquad \pi(\phi_{W,f})=W,
\]
is well-defined and $U(N)$-equivariant.
\item[(iii)] For fixed $W$, the fiber $\pi^{-1}(W)$ is naturally a torsor for
$PSU(2)=SU(2)/\{\pm I\}\cong SO(3)$. Equivalently, $\pi$ exhibits
$\Emb_{\mathrm{2lvl}}(SU(2),U(N))$ as a principal $PSU(2)$-bundle over $\Gr_2(\C^N)$.
\end{enumerate}
\end{proposition}

\begin{proof}
(i) Any two-planes $W,W'$ are related by some $U\in U(N)$. Conjugation by $U$ transports $SU(N)[W]$ to
$SU(N)[W']$, hence carries $\phi_{W,f}$ to an embedding of the form $\phi_{W',f'}$ for some frame
$f'$. Thus conjugation is transitive.

(ii) By Proposition~\ref{prop:gauge_invariance_SU2}, the subgroup $\phi_{W,f}(SU(2))$ depends only on
$W$, so $\pi$ is well-defined. Equivariance is immediate.

(iii) Fix $W$ and two frames $f,f'$ over $W$. Then $f'=fU$ for a unique $U\in U(2)$, and by the proof
of Proposition~\ref{prop:gauge_invariance_SU2} we obtain
$\phi_{W,f'}=\phi_{W,f}\circ \Ad_{S_0}$ for some $S_0\in SU(2)$ well-defined modulo $\{\pm I\}$.
Hence the effective gauge group is $PSU(2)$ and $\pi^{-1}(W)$ is a $PSU(2)$-torsor.
\end{proof}


%------------------------------------------------------------
\subsection{Descriptor form: separating the logical support from the internal action}
\label{subsec:descriptor_form}
%------------------------------------------------------------

Two-level gates admit a convenient \emph{descriptor} that separates \emph{where} the gate acts (the
support two-plane) from \emph{what} action is performed on that plane (the internal $SU(2)$ element),
while keeping track of the intrinsic gauge redundancy.

\paragraph{Support as a projector (gauge-invariant data).}
Given $W\in\Gr_2(\C^N)$, let $P_W$ denote the orthogonal projector onto $W$. The map $W\mapsto P_W$
identifies $\Gr_2(\C^N)$ with the set of rank-$2$ orthogonal projectors in $\End(\C^N)$; thus $P_W$ is
an intrinsic representative of the logical-qubit choice.

\paragraph{Internal action and gauge.}
Fix $W$ and choose a frame $f:\C^2\to W$. Any two-level special unitary supported on $W$ can be
written as
\begin{equation}\label{eq:descriptor_group_level_SU2}
U_{W,f}(S)\;:=\;\phi_{W,f}(S)\;=\; fSf^{-1}\oplus I_{W^\perp},\qquad S\in SU(2).
\end{equation}
If the frame is changed to $f'=fU$ with $U\in U(2)$, then $U_{W,f'}(S)=U_{W,f}(S_0SS_0^{-1})$ for a
suitable $S_0\in SU(2)$, so the internal label is defined only up to conjugation, i.e.\ up to the
$PSU(2)$ gauge.

\paragraph{Lie-algebra descriptor.}
For geometric and variational constructions it is convenient to work at Lie-algebra level. Given
$A\in\mathfrak{su}(2)$, define the corresponding two-level generator on $\C^N$ by
\begin{equation}\label{eq:descriptor_lie_level_SU2}
X_{W,f}(A)\;:=\; fAf^{-1}\oplus 0_{W^\perp}\ \in\ \mathfrak{u}(N),
\qquad
U_{W,f}\big(\exp(A)\big)=\exp\!\big(X_{W,f}(A)\big).
\end{equation}

\begin{lemma}[Hilbert--Schmidt isometry for two-level generators]
\label{lem:HS_norm_two_level_generator}
Let $W\in\Gr_2(\C^N)$, let $f:\C^2\to W$ be a unitary frame, and let $A,B\in\mathfrak{su}(2)$.
Define $X_{W,f}(A):=fAf^{-1}\oplus 0_{W^\perp}$ as in \eqref{eq:descriptor_lie_level_SU2}. Then, with
$\langle X,Y\rangle_{\mathrm{hs}}=\tfrac12\Tr(X^\dagger Y)$ on both $\mathfrak{su}(2)$ and
$\mathfrak{u}(N)$,
\begin{equation}\label{eq:HS_norm_preservation}
\|X_{W,f}(A)\|_{\mathrm{hs}}=\|A\|_{\mathrm{hs}}
\qquad\text{and}\qquad
\langle X_{W,f}(A),X_{W,f}(B)\rangle_{\mathrm{hs}}=\langle A,B\rangle_{\mathrm{hs}}.
\end{equation}
In particular, these quantities are independent of the choice of frame $f$.
\end{lemma}

\begin{proof}
Since $f$ is unitary, conjugation preserves adjoints and traces. Writing
$X_{W,f}(A)=fAf^{-1}\oplus 0$, we have
\[
\|X_{W,f}(A)\|_{\mathrm{hs}}^2
=\frac12\Tr\!\big((fA^\dagger f^{-1})(fAf^{-1})\big)
=\frac12\Tr(A^\dagger A)
=\|A\|_{\mathrm{hs}}^2,
\]
and similarly for the inner product. Frame-independence follows.
\end{proof}

\begin{corollary}[Energy reduction for two-level paths]\label{cor:energy_reduction_two_level}
Let $A:[0,1]\to\mathfrak{su}(2)$ be measurable and consider a curve $U:[0,1]\to U(N)$ with body
velocity $X_{W,f}(A(t))$. Then
\[
\int_0^1\|X_{W,f}(A(t))\|_{\mathrm{hs}}^2\,dt
=\int_0^1\|A(t)\|_{\mathrm{hs}}^2\,dt,
\]
so the Hilbert--Schmidt metric cost of a two-level evolution is exactly the cost of its
$2\times 2$ internal generator.
\end{corollary}

\begin{remark}[A gauge-invariant descriptor pair]\label{rem:gauge_invariant_descriptor_pair}
A two-level special-unitary gate supported on $W$ can be described by the pair $(P_W,[S])$, where
$P_W$ is the rank-$2$ projector onto $W$ and $[S]$ is the $PSU(2)$-class of $S$ (i.e.\ $S$ modulo the
conjugation gauge). Thus $P_W$ specifies the logical support, while $[S]$ specifies the internal
action up to the frame choice.
\end{remark}

\paragraph{Phase conventions.}
If one wishes to incorporate arbitrary phases on the support plane $W$ (i.e.\ to work with $U(N)[W]$
rather than $SU(N)[W]$), one may extend the $SU(2)$ action by an additional commuting $U(1)$ phase on
$W$. Equivalently, one may treat the $SU(2)$ two-level part as the elementary non-abelian primitive
and compile the remaining abelian degrees of freedom (global phase and/or diagonal factors)
separately, as described in the discretization/compilation step.


%------------------------------------------------------------
\subsection{Relation to the embedding dictionary and to two-level unitaries}
\label{subsec:relation_to_dict}
%------------------------------------------------------------

Two-level embeddings yield a Grassmannian-indexed subset of the full embedding-elementary dictionary:
\[
\begin{aligned}
\mathcal{G}^{SU}_{\mathrm{2lvl}}(n)
&:=\bigcup_{W\in\Gr_2(\C^N)} SU(N)[W],\\
\mathcal{G}^{U}_{\mathrm{2lvl}}(n)
&:=\bigcup_{W\in\Gr_2(\C^N)} U(N)[W].
\end{aligned}
\]
Elements of $SU(N)[W]$ are precisely the \emph{special-unitary two-level unitaries}: they act as an
arbitrary $SU(2)$ on $W$ and as the identity on $W^\perp$. These operations are the basic primitives
in classical two-level (Givens-type) synthesis and will be the core mechanism behind the universality
proofs in Section~\ref{sec:universality}.

\begin{proposition}[Two-level gates are embedding-elementary]\label{prop:two_level_in_Gelem_SU2}
For every $W\in\Gr_2(\C^N)$ one has $SU(N)[W]\subset \mathcal{G}_{\mathrm{elem}}^{SU}(n)$. Consequently,
\[
\mathcal{G}_{\mathrm{2lvl}}^{SU}(n)\subset \mathcal{G}_{\mathrm{elem}}^{SU}(n).
\]
\end{proposition}

\begin{proof}
Fix $W$ and choose a frame $f:\C^2\to W$. Then $SU(N)[W]=\phi_{W,f}(SU(2))$ and
$\phi_{W,f}\in\Emb(SU(2),U(N))$, hence every element of $SU(N)[W]$ lies in
$\bigcup_{\phi\in\Emb(SU(2),U(N))}\phi(SU(2))=\mathcal{G}_{\mathrm{elem}}^{SU}(n)$.
\end{proof}


%------------------------------------------------------------
\subsection{Comparison with tensor-factor locality}
\label{subsec:comparison_tensor_locality}
%------------------------------------------------------------

The Grassmannian model is intrinsic on $\C^N$ and does not presuppose a tensor factorization.
Nevertheless, once $\mathbb{H}_N=(\C^2)^{\otimes n}$ is fixed, conventional one-qubit gates appear as
a distinguished embedding-elementary family, although typically not as two-level operations.

\begin{proposition}[Tensor-local one-qubit gates are embedding-elementary (but not two-level)]
\label{prop:local_in_Gelem_not_2lvl}
Let $w_j^{(n)}:U(2)\hookrightarrow U(N)$ be the canonical tensor placements from
Section~\ref{sec:preliminaries}. Then $w_j^{(n)}(SU(2))\subset \mathcal{G}_{\mathrm{elem}}^{SU}(n)$
for each $j$, and hence
\[
\bigcup_{j=1}^n w_j^{(n)}(SU(2))\ \subset\ \mathcal{G}_{\mathrm{elem}}^{SU}(n).
\]
Moreover, for $n\ge 2$ the subgroup $w_j^{(n)}(SU(2))$ is not of two-level form $SU(N)[W]$: it acts
nontrivially on a $2^{n-1}$-fold orthogonal direct sum of $2$-planes (the subspaces spanned by
$\ket{0,b}$ and $\ket{1,b}$ with $b\in\{0,1\}^{n-1}$), rather than on a single two-plane with
pointwise-fixed orthogonal complement.
\end{proposition}

\begin{proof}
By definition, $w_j^{(n)}$ is an embedding $U(2)\hookrightarrow U(N)$, hence its restriction to
$SU(2)$ lies in $\Emb(SU(2),U(N))$, so $w_j^{(n)}(SU(2))\subset \mathcal{G}_{\mathrm{elem}}^{SU}(n)$.

For the final claim, observe that $w_j^{(n)}(SU(2))$ preserves the decomposition
$\C^N\cong\bigoplus_{b\in\{0,1\}^{n-1}} \C^2_{b}$ into $2^{n-1}$ orthogonal two-planes, acting as the
\emph{same} $SU(2)$ on each plane $\C^2_b$. In contrast, a two-level subgroup $SU(N)[W]$ fixes
$W^\perp$ pointwise and acts nontrivially only on a single two-plane $W$. Thus $w_j^{(n)}(SU(2))$ is
not of the form $SU(N)[W]$ when $n\ge 2$.
\end{proof}

\begin{remark}[Why the Grassmannian model is the canonical two-level notion]
Tensor locality selects $n$ preferred embeddings of $SU(2)$ (the tensor factors), but these lie in a
different representation type than the two-level stratum: they act as multiple copies of the defining
module rather than as $V_1\oplus V_0^{\oplus(N-2)}$. The Grassmannian two-level model instead ranges
over \emph{all} complex two-planes and isolates exactly the sector underlying classical two-level
(Givens-type) synthesis and the corresponding universality mechanism for $SU(N)$.
\end{remark}

\paragraph{Bridge to universality and compilation.}
The key point for the remainder of the paper is that two-level gates provide both (i) an intrinsic,
Grassmannian-indexed primitive vocabulary and (ii) the classical two-level factors used in
constructive synthesis. In Section~\ref{sec:universality} we prove that the subgroup generated by
$\bigcup_{W}SU(N)[W]$ is $SU(N)$, and in Section~\ref{sec:SK} we discretize these $SU(2)$ motions and
lift the resulting words to $U(N)$ with explicit error control.


%============================================================
\section{Phase-free and full universality from two-level dictionaries}
\label{sec:universality}
%============================================================

\noindent
We prove that the two-level sector already yields universality, both in the phase-free sense (generation of $SU(N)$) and in the full
unitary sense (generation of $U(N)$). Throughout, fix $n\in\N$ and set $N:=2^n$. Recall that for each $2$-plane
$W\in\Gr_2(\C^N)$ we have the two-level subgroups $U(N)[W]\le U(N)$ and $SU(N)[W]\le SU(N)$
(Definition~\ref{def:two_level_subgroup}). We consider the associated two-level dictionaries
\begin{equation}\label{eq:G2lvl_def_SU_U_concise}
\mathcal{G}^{SU}_{\mathrm{2lvl}}(n):=\bigcup_{W\in\Gr_2(\C^N)} SU(N)[W]\subset SU(N),
\qquad
\mathcal{G}^{U}_{\mathrm{2lvl}}(n):=\bigcup_{W\in\Gr_2(\C^N)} U(N)[W]\subset U(N).
\end{equation}
Since two-level subgroups are realized as images of (framed) two-level embeddings, we have
$\mathcal{G}^{SU}_{\mathrm{2lvl}}(n)\subset \mathcal{G}^{SU}_{\mathrm{elem}}(n)$ and
$\mathcal{G}^{U}_{\mathrm{2lvl}}(n)\subset \mathcal{G}^{U}_{\mathrm{elem}}(n)$
(cf.\ Sections~\ref{subsec:two_level_embeddings}--\ref{subsec:relation_to_dict}). Hence universality at the two-level level implies
universality of the full embedding-elementary dictionaries.

\medskip
\noindent
The argument has two inputs: (i) a QR/Givens factorization of arbitrary unitaries into coordinate two-level factors and a diagonal
remainder, and (ii) the fact that diagonal tori are generated by two-level phase rotations. We treat $SU(N)$ and $U(N)$ in parallel.

%------------------------------------------------------------
\subsection{Diagonal tori from two-level phase rotations}
\label{subsec:diagonal_tori_two_level}
%------------------------------------------------------------

Let $T\le U(N)$ be the diagonal torus and $T_0:=T\cap SU(N)$.

\begin{lemma}[The special-unitary torus is generated by $SU(2)$ two-level phases]
\label{lem:T0_generated_by_two_level_phases_concise}
For each $j=2,\dots,N$ define
\[
H_{1j}:=\frac{i}{2}\big(E_{11}-E_{jj}\big)\in\mathfrak{su}(N),
\qquad
\gamma_{1j}(t):=\exp(tH_{1j})\in SU(N).
\]
Then $\gamma_{1j}(t)\in SU(N)[W_{1j}]$ for $W_{1j}=\mathrm{span}\{e_1,e_j\}$, and
\[
T_0=\Big\langle\,\gamma_{12}(\R),\dots,\gamma_{1N}(\R)\,\Big\rangle.
\]
Consequently, every $D\in T$ admits a factorization $D=e^{i\theta}D_0$ with $\theta\in\R$ and
$D_0\in\big\langle\,SU(N)[W_{12}]\allowbreak\cup\cdots\allowbreak\cup SU(N)[W_{1N}]\,\big\rangle$.
\end{lemma}

\begin{proof}
The support statement follows from the explicit diagonal action of $H_{1j}$, hence of $\gamma_{1j}(t)$.
Moreover, $\mathfrak{t}_0=\Lie(T_0)$ is spanned by $\{\,i(E_{11}-E_{jj})\,\}_{j=2}^N$, equivalently by $\{H_{1j}\}_{j=2}^N$, and the
$H_{1j}$ commute. Therefore $\exp(\mathfrak{t}_0)$ lies in the subgroup generated by the one-parameter subgroups
$\gamma_{1j}(\R)$. Since $T_0$ is a compact torus, $\exp:\mathfrak{t}_0\to T_0$ is surjective
\cite[Ch.~7]{Hall2015} or \cite[Ch.~IV]{Knapp2002}, giving the claim. The final statement is the decomposition of $T$ into global phase
and special-unitary part.
\end{proof}

\begin{lemma}[The full diagonal torus is generated by $U(2)$ two-level phases]
\label{lem:T_generated_by_U2_two_level_phases_concise}
One has $T\subset \langle \mathcal{G}^{U}_{\mathrm{2lvl}}(n)\rangle$.
\end{lemma}

\begin{proof}
For each $j$ and $t$ the diagonal unitary $\eta_j(t):=\diag(1,\dots,e^{it},\dots,1)$ is supported on $W_{1j}$ (for $j\ge 2$) as a
two-level $U(2)$ phase rotation, hence belongs to $\mathcal{G}^{U}_{\mathrm{2lvl}}(n)$. Since $T$ is generated by these one-parameter
subgroups, the result follows.
\end{proof}

%------------------------------------------------------------
\subsection{Coordinate two-level factorization with diagonal remainder}
\label{subsec:two_level_factorization}
%------------------------------------------------------------

\begin{theorem}[Unitary QR/Givens factorization into coordinate two-level gates]
\label{thm:two_level_factorization_concise}
Fix an orthonormal basis $\{e_1,\dots,e_N\}$ of $\C^N$. Then every $U\in U(N)$ admits a factorization
\begin{equation}\label{eq:U_two_level_factorization_concise}
U=\Big(\prod_{k=1}^{K} T_k\Big)\,D,
\qquad
T_k\in U(N)[W_{p_k,q_k}],\ \ D\in T,
\end{equation}
where each $W_{p_k,q_k}=\mathrm{span}\{e_{p_k},e_{q_k}\}$ is a coordinate two-plane and $K\le \frac{N(N-1)}{2}$.
Moreover, the factors can be chosen by standard Givens elimination to annihilate prescribed subdiagonal entries.
\end{theorem}

\begin{proof}
This is the standard unitary QR procedure via Givens rotations: successive left-multiplication by coordinate two-level unitaries
eliminates subdiagonal entries, yielding an upper-triangular unitary matrix, hence a diagonal matrix $D\in T$.
See, e.g., \cite{GolubVanLoan2013,Higham2008,HornJohnson2013}.
\end{proof}

%------------------------------------------------------------
\subsection{Universality: phase-free and full}
\label{subsec:universality_phase_free_full}
%------------------------------------------------------------

\begin{theorem}[Phase-free universality of $\mathcal{G}^{SU}_{\mathrm{2lvl}}(n)$]
\label{thm:SU_two_level_universal_concise}
For $N=2^n$,
\[
\big\langle \mathcal{G}^{SU}_{\mathrm{2lvl}}(n)\big\rangle = SU(N).
\]
\end{theorem}

\begin{proof}
Let $U\in SU(N)$. By Theorem~\ref{thm:two_level_factorization_concise},
$U=(\prod_{k=1}^K T_k)D$ with $T_k\in U(N)[W_{p_k,q_k}]$ and $D\in T$. Since $\det U=1$, necessarily $D\in T_0$.

For each $k$, write $\det(T_k|_{W_{p_k,q_k}})=e^{i\vartheta_k}$ and set $s_k:=e^{-i\vartheta_k/2}$.
Then $T_k=(S_k\oplus I)\Delta_k$ where $S_k\oplus I\in SU(N)[W_{p_k,q_k}]\subset\mathcal{G}^{SU}_{\mathrm{2lvl}}(n)$ and $\Delta_k\in T$
is diagonal. Absorbing $\prod_k\Delta_k$ into the trailing diagonal term yields a new diagonal factor $D_0\in T$ with $\det(D_0)=1$,
hence $D_0\in T_0$. By Lemma~\ref{lem:T0_generated_by_two_level_phases_concise}, $D_0$ lies in the subgroup generated by two-level
$SU(2)$ phase rotations. Therefore $U$ is a product of elements of $\mathcal{G}^{SU}_{\mathrm{2lvl}}(n)$.
\end{proof}

\begin{theorem}[Full universality of $\mathcal{G}^{U}_{\mathrm{2lvl}}(n)$]
\label{thm:U_two_level_universal_concise}
For $N=2^n$,
\[
\big\langle \mathcal{G}^{U}_{\mathrm{2lvl}}(n)\big\rangle = U(N).
\]
\end{theorem}

\begin{proof}
Let $U\in U(N)$. Theorem~\ref{thm:two_level_factorization_concise} gives $U=(\prod_{k=1}^K T_k)D$ with each
$T_k\in \mathcal{G}^{U}_{\mathrm{2lvl}}(n)$ and $D\in T$. By Lemma~\ref{lem:T_generated_by_U2_two_level_phases_concise},
$D\in\langle \mathcal{G}^{U}_{\mathrm{2lvl}}(n)\rangle$. Hence $U\in\langle \mathcal{G}^{U}_{\mathrm{2lvl}}(n)\rangle$.
\end{proof}

\begin{corollary}[Universality of the embedding-elementary dictionaries]
\label{cor:universality_elem_from_2lvl_concise}
Let $N=2^n$. Then
\[
\begin{aligned}
\big\langle \mathcal{G}^{SU}_{\mathrm{2lvl}}(n)\big\rangle &= SU(N),\\
\big\langle \mathcal{G}^{U}_{\mathrm{2lvl}}(n)\big\rangle &= U(N).
\end{aligned}
\]
\end{corollary}

\begin{proof}
By $\mathcal{G}^{SU}_{\mathrm{2lvl}}(n)\subset \mathcal{G}^{SU}_{\mathrm{elem}}(n)\subset SU(N)$ and
$\mathcal{G}^{U}_{\mathrm{2lvl}}(n)\subset \mathcal{G}^{U}_{\mathrm{elem}}(n)\subset U(N)$, combine
Theorems~\ref{thm:SU_two_level_universal_concise} and~\ref{thm:U_two_level_universal_concise}.
\end{proof}

%------------------------------------------------------------
\subsection{Contrast: strict locality alone is not universal}
\label{subsec:local_not_universal}
%------------------------------------------------------------

Recall $\mathcal{G}_{\mathrm{loc}}(n)=\bigcup_{j=1}^n w_j^{(n)}(U(2))\subset U(N)$.

\begin{proposition}[Local one-qubit gates generate only tensor-product unitaries]
\label{prop:Gloc_generated_group_concise}
For $n\ge 1$,
\[
\big\langle \mathcal{G}_{\mathrm{loc}}(n)\big\rangle
=
\{\,U_1\otimes\cdots\otimes U_n:\ U_j\in U(2)\,\}
\ \cong\ U(2)^n.
\]
\end{proposition}

\begin{proof}
The subgroups $w_j^{(n)}(U(2))$ act on different tensor factors and therefore commute. The subgroup generated by their union is their
direct product, i.e.\ the set of tensor-product unitaries.
\end{proof}

\begin{corollary}[Non-universality of $\mathcal{G}_{\mathrm{loc}}(n)$ for $n\ge2$]
\label{cor:Gloc_not_universal_concise}
If $n\ge2$, then $\langle \mathcal{G}_{\mathrm{loc}}(n)\rangle\subsetneq U(2^n)$.
\end{corollary}

\begin{proof}
$\langle \mathcal{G}_{\mathrm{loc}}(n)\rangle\cong U(2)^n$ contains no entangling unitaries, hence is a proper subgroup of $U(2^n)$
for $n\ge2$.
\end{proof}

\begin{remark}[Interface with discretization]
\label{rem:universality_to_discretization_concise}
The factorization \eqref{eq:U_two_level_factorization_concise} reduces synthesis in $U(N)$ to the synthesis of a sequence of $U(2)$
blocks supported on two-planes. Fixing a finite universal set $G_2\subset SU(2)$, each $SU(2)$ block can be approximated by words over
$G_2$ (e.g.\ Solovay--Kitaev), and the resulting words lift to $U(N)$ through the corresponding two-level embeddings. This is the
starting point for the finite-alphabet compilation framework of Section~\ref{sec:SK}.
\end{remark}

%============================================================
\subsection{Standard controlled and permutation gates as embedded $U(2)$ elements}
\label{subsec:standard_gates_as_embeddings}
%============================================================

\noindent
As a consistency check, common circuit primitives arise transparently as embedded $U(2)$ elements: they act as the identity on a large
orthogonal subspace and apply a $2\times2$ block on a distinguished two-plane selected by a classical predicate on basis states. This is
exactly the structure of a faithful unitary representation of $U(2)$ on $\C^N$, hence an element of $\Emb(U(2),U(N))$
(cf.\ Proposition~\ref{prop:embeddings_as_faithful_representations_G}).


%============================================================
\section{Discretization and Compilation: Two-Level Factorizations and the Solovay--Kitaev Paradigm}
\label{sec:SK}
%============================================================

\noindent
The embedding-elementary dictionary $\mathcal{G}_{\mathrm{elem}}(n)$ is a \emph{continuous} generating family in $U(N)$.
To obtain an implementable gate model, one must pass to a \emph{finite} alphabet and provide a compilation procedure with explicit
accuracy--length guarantees. We record a modular route consistent with the Grassmannian two-level viewpoint:
\begin{enumerate}
\item[(1)] (\emph{Exact reduction to two-level factors}) every $U\in U(N)$ admits an \emph{exact} factorization into coordinate two-level
unitaries via unitary QR/Givens eliminations;
\item[(2)] (\emph{Local synthesis problem}) approximating each two-level factor reduces compilation in $U(N)$ to repeated approximation
problems in $U(2)$ (or, after phase normalization, in $SU(2)$);
\item[(3)] (\emph{Finite alphabet on $SU(2)$}) the Solovay--Kitaev theorem yields polylogarithmic-length approximation of $SU(2)$ elements
by words over a fixed finite universal set;
\item[(4)] (\emph{Lift and control the global error}) lifting these words through the corresponding two-level embeddings $\phi_{p,q}$
produces an $\varepsilon$-approximation in $U(N)$ with explicit operator-norm error control, while global phase and diagonal remainders
are handled separately as needed.
\end{enumerate}

%============================================================
\subsection{From a finite $SU(2)$ alphabet to a finite elementary alphabet in $U(N)$}
\label{subsec:finite_alphabet}
%============================================================

Let $G_2\subset SU(2)$ be a finite set. We say that $G_2$ is \emph{universal for $SU(2)$} if $\langle G_2\rangle$ is dense in $SU(2)$.
Given such a set, we define the induced finite elementary alphabet in $U(N)$ by
\begin{equation}\label{eq:finite_embedded_alphabet}
\mathcal{A}_{\mathrm{elem}}(n;G_2)
\;:=\;
\bigcup_{1\le p<q\le N}\phi_{p,q}(G_2)\ \subset\ U(N).
\end{equation}
Thus $\mathcal{A}_{\mathrm{elem}}(n;G_2)$ consists of \emph{discrete two-level gates} acting by one of finitely many $SU(2)$ primitives
on a chosen coordinate two-plane and trivially on its orthogonal complement.

\begin{remark}[Determinants and diagonal phases]\label{rem:det_bookkeeping}
Each $V\in U(2)$ admits a factorization $V=e^{i\theta}\widetilde V$ with $\widetilde V\in SU(2)$ (take $e^{i\theta}=\sqrt{\det V}$).
In the exact synthesis \eqref{eq:U_two_level_factorization_concise}, the scalar phases $e^{i\theta_k}$ may be absorbed into the trailing diagonal
unitary $D$. Accordingly, the approximation results below are stated for $SU(2)$ factors and extend to $U(2)$ with this bookkeeping.
\end{remark}

%============================================================
\subsection{Solovay--Kitaev approximation on $SU(2)$}
\label{subsec:SK_statement}
%============================================================

We write $\|\cdot\|$ for the operator norm (spectral norm) induced by the Euclidean norm on $\C^m$.

\begin{theorem}[Solovay--Kitaev theorem for $SU(2)$]\label{thm:SK}
Let $G_2\subset SU(2)$ be a finite set such that the subgroup $\langle G_2\rangle$ is dense in $SU(2)$.
Then there exist constants $c,C>0$, depending only on $G_2$, and a deterministic classical algorithm with the following property:
for every $V\in SU(2)$ and every $\varepsilon\in(0,1)$ the algorithm produces a word $w=w(V,\varepsilon)$ over
$G_2\cup G_2^{-1}$ such that
\[
\|V-w\|\le \varepsilon,
\qquad
|w|\le C\,\log^{c}\!\Big(\frac{1}{\varepsilon}\Big).
\]
Moreover, the running time is polynomial in $\log(1/\varepsilon)$.
\end{theorem}

\begin{remark}
In what follows we only use the existence of a polylogarithmic length bound in $\varepsilon^{-1}$ and the fact that the word can be
found effectively. See \cite{Kitaev2002,DawsonNielsen2006} for constructive versions and explicit exponents.
\end{remark}

While we only use Solovay--Kitaev as a modular finite-alphabet interface for approximating $SU(2)$ primitives,
it is worth noting that the asymptotic word-length exponent in the general Solovay--Kitaev theorem has been
substantially improved very recently: Kuperberg breaks the classical ``cubic'' barrier and obtains word length
$\ell = O(n^{\alpha+\delta})$ with $\alpha> \log_{\varphi}(2)\approx 1.44042$ (for inverse-closed finite gate sets
densely generating $SU(d)$, and more generally connected semisimple Lie groups), together with a near-linear
bound for the corresponding compressed-word length; see~\cite{Kuperberg2025BreakingCubicSK}.
These refinements can be plugged into our pipeline verbatim whenever one wants sharper asymptotic gate counts
for the $SU(2)$ approximation subroutine.

%============================================================
\subsection{Lifting Solovay--Kitaev words to $U(N)$ and global error control}
\label{subsec:lifting_and_error}
%============================================================

The following elementary lemmas isolate the analytic inputs needed to propagate local approximation errors
to a global operator-norm bound. Throughout, $\|\cdot\|$ denotes the operator norm on $\mathbb{C}^N$; we repeatedly use its unitary
invariance $\|UXV\|=\|X\|$ for $U,V\in U(N)$ and its submultiplicativity $\|XY\|\le \|X\|\,\|Y\|$ (see, e.g., \cite{HornJohnson2013}).

\begin{lemma}[Operator-norm isometry for coordinate two-level embeddings]
\label{lem:two_level_isometry}
Let $1\le p<q\le N$ and let $\phi_{p,q}:U(2)\hookrightarrow U(N)$ be the coordinate two-level embedding supported on
$W_{p,q}=\mathrm{span}\{e_p,e_q\}$. Then for all $V,W\in U(2)$,
\[
\big\|\phi_{p,q}(V)-\phi_{p,q}(W)\big\|=\|V-W\|.
\]
In particular, $\phi_{p,q}$ is an isometric embedding with respect to the operator norm.
\end{lemma}

\begin{proof}
By construction, there exists a unitary $Q\in U(N)$ such that, with respect to the orthogonal decomposition
$\mathbb{C}^N=W_{p,q}\oplus W_{p,q}^\perp$, one has
\[
\phi_{p,q}(V)=Q^{-1}\,\mathrm{diag}(V,I_{N-2})\,Q,
\qquad
\phi_{p,q}(W)=Q^{-1}\,\mathrm{diag}(W,I_{N-2})\,Q.
\]
Hence, by unitary invariance of the operator norm,
\[
\|\phi_{p,q}(V)-\phi_{p,q}(W)\|
=
\big\|\,\mathrm{diag}(V-W,0)\,\big\|.
\]
For a block-diagonal operator the operator norm equals the maximum of the norms of the diagonal blocks, so
$\big\|\,\mathrm{diag}(V-W,0)\,\big\|=\|V-W\|$.
\end{proof}

\begin{lemma}[Telescoping bound for products]
\label{lem:error_accumulation}
Let $A_1,\dots,A_K,B_1,\dots,B_K\in U(N)$. Then
\[
\Big\|\prod_{k=1}^K A_k-\prod_{k=1}^K B_k\Big\|
\ \le\
\sum_{k=1}^K \|A_k-B_k\|.
\]
\end{lemma}

\begin{proof}
Using the telescoping identity
\[
\prod_{k=1}^K A_k-\prod_{k=1}^K B_k
=
\sum_{j=1}^K
\Big(\prod_{k=1}^{j-1} B_k\Big)\,(A_j-B_j)\,\Big(\prod_{k=j+1}^{K} A_k\Big),
\]
and taking norms yields, by the triangle inequality and submultiplicativity,
\[
\Big\|\prod_{k=1}^K A_k-\prod_{k=1}^K B_k\Big\|
\le
\sum_{j=1}^K
\Big\|\prod_{k=1}^{j-1} B_k\Big\|\,\|A_j-B_j\|\,\Big\|\prod_{k=j+1}^{K} A_k\Big\|.
\]
If all factors are unitary then each partial product is unitary and hence has norm $1$, giving the claim.
\end{proof}

We now combine the exact two-level factorization (Theorem~\ref{thm:two_level_factorization_concise}) with Solovay--Kitaev approximation
(Theorem~\ref{thm:SK}).

\begin{theorem}[Compilation via two-level factorization and Solovay--Kitaev]
\label{thm:UN_compilation_SK}
Fix $N=2^n$ and let $G_2\subset SU(2)$ be a finite universal set.
Let $U\in U(N)$ and $\varepsilon\in(0,1)$. Then there exist:
\begin{itemize}
\item a word $W(U,\varepsilon)$ over the finite alphabet
\[
\mathcal{A}_{\mathrm{elem}}(n;G_2)\ \cup\ \mathcal{A}_{\mathrm{elem}}(n;G_2)^{-1},
\]
\item and a diagonal unitary $D\in U(N)$,
\end{itemize}
such that
\begin{equation}\label{eq:global_approx}
\big\|U - W(U,\varepsilon)\,D\big\|\ \le\ \varepsilon.
\end{equation}
Moreover, one may choose $W(U,\varepsilon)$ with length bounded by
\begin{equation}\label{eq:length_bound_global}
|W(U,\varepsilon)|\ \le\ K\cdot C\,\log^{\,c}\!\Big(\frac{K}{\varepsilon}\Big),
\qquad
K\le \frac{N(N-1)}{2},
\end{equation}
where $c,C$ are the Solovay--Kitaev constants from Theorem~\ref{thm:SK}.
If one wishes to eliminate the residual diagonal factor, it can be compiled (up to global phase) using two-level phase rotations as in
Section~\ref{subsec:compile_diagonal_optional}.
\end{theorem}

\begin{proof}
By Theorem~\ref{thm:two_level_factorization_concise} there exist indices $(p_k,q_k)$, factors $V_k\in U(2)$, and a diagonal unitary $D\in U(N)$
such that
\[
U=\Big(\prod_{k=1}^K \phi_{p_k,q_k}(V_k)\Big)\,D,
\qquad
K\le \frac{N(N-1)}{2}.
\]
Using Remark~\ref{rem:det_bookkeeping}, write $V_k=e^{i\theta_k}\widetilde V_k$ with $\widetilde V_k\in SU(2)$ and absorb the phases
into $D$, so that we may assume $V_k\in SU(2)$ without loss of generality.

Set $\delta:=\varepsilon/K$. For each $k$, Theorem~\ref{thm:SK} yields a word $w_k$ over $G_2\cup G_2^{-1}$ such that
\[
\|V_k-w_k\|\le \delta,
\qquad
|w_k|\le C\,\log^{c}\!\Big(\frac{1}{\delta}\Big)=C\,\log^{c}\!\Big(\frac{K}{\varepsilon}\Big).
\]
Define the lifted word in $U(N)$ by
\[
W(U,\varepsilon):=\prod_{k=1}^K \phi_{p_k,q_k}(w_k),
\]
where $\phi_{p_k,q_k}(w_k)$ denotes the word obtained by embedding each letter of $w_k$ through $\phi_{p_k,q_k}$.
By Lemma~\ref{lem:two_level_isometry},
\[
\big\|\phi_{p_k,q_k}(V_k)-\phi_{p_k,q_k}(w_k)\big\|
=
\|V_k-w_k\|
\le \delta.
\]
Applying Lemma~\ref{lem:error_accumulation} gives
\[
\Big\|\prod_{k=1}^K\phi_{p_k,q_k}(V_k)-\prod_{k=1}^K\phi_{p_k,q_k}(w_k)\Big\|
\le
\sum_{k=1}^K \delta
=
\varepsilon.
\]
Right-multiplication by $D$ preserves the operator norm (unitary invariance), hence \eqref{eq:global_approx} follows.
Finally,
\[
|W(U,\varepsilon)|
=\sum_{k=1}^K |w_k|
\le
K\cdot C\,\log^{c}\!\Big(\frac{K}{\varepsilon}\Big),
\]
which is \eqref{eq:length_bound_global}.
\end{proof}

\begin{remark}[Worst-case scaling in $n$]\label{rem:worst_case_scaling}
Since $K\le N(N-1)/2=\Theta(N^2)=\Theta(4^n)$, the bound \eqref{eq:length_bound_global} yields worst-case word lengths that scale
exponentially in $n$ for generic targets in $U(2^n)$. This reflects the parameter complexity of $U(2^n)$ and should be interpreted
as providing a representation-invariant and modular discretization mechanism rather than an efficiency guarantee for arbitrary
instances.
\end{remark}

%============================================================
\subsection{Compiling the diagonal factor}
\label{subsec:compile_diagonal_optional}
%============================================================

In the QR/Givens synthesis of Theorem~\ref{thm:UN_compilation_SK} the target unitary $U\in U(N)$ is produced up to a trailing diagonal
unitary $D$. If one wishes to obtain a ``pure word'' over embedded $SU(2)$ generators (up to an unavoidable global phase), then $D$ can
itself be synthesized by two-level phase rotations supported on coordinate planes. We record a convenient formulation.


\begin{corollary}[Pure-word compilation up to global phase]
\label{cor:pure_word_compilation}
Under the assumptions of Theorem~\ref{thm:UN_compilation_SK}, one may absorb the diagonal factor into the word up to a global phase:
there exist $\theta\in\mathbb{R}$ and a word $\widetilde W(U,\varepsilon)$ over
$\mathcal{A}_{\mathrm{elem}}(n;G_2)\cup\mathcal{A}_{\mathrm{elem}}(n;G_2)^{-1}$ such that
\[
\big\|U-e^{i\theta}\widetilde W(U,\varepsilon)\big\|\le \varepsilon.
\]
\end{corollary}

\begin{proof}
Apply Theorem~\ref{thm:UN_compilation_SK} to obtain $U=W(U,\varepsilon)\,D+E$ with $\|E\|\le\varepsilon$.
Decompose $D=e^{i\theta}D_0$ as in Lemma~\ref{lem:T0_generated_by_two_level_phases_concise} and write $D_0$ as a finite product of
elements from $\phi_{1j}(SU(2))$ (hence as a word over the corresponding embedded alphabet).
Let $\widetilde W(U,\varepsilon):=W(U,\varepsilon)\cdot(\text{word for }D_0)$. Then $e^{i\theta}\widetilde W=W D$ by construction, so
$U-e^{i\theta}\widetilde W=E$ and the norm bound follows.
\end{proof}

%============================================================
\subsection{Algorithmic summary and interpretation}
\label{subsec:algorithmic_summary}
%============================================================

We summarize the construction of Theorem~\ref{thm:UN_compilation_SK} as an explicit compilation procedure.

\medskip
\noindent\textbf{Compilation procedure (coordinate two-level factorization + Solovay--Kitaev).}
Given $U\in U(N)$, $\varepsilon\in(0,1)$, and a finite universal set $G_2\subset SU(2)$, the following steps produce
a word $W$ over $\mathcal{A}_{\mathrm{elem}}(n;G_2)\cup\mathcal{A}_{\mathrm{elem}}(n;G_2)^{-1}$ and a diagonal $D\in U(N)$ such that
$\|U-WD\|\le \varepsilon$.

\begin{quote}
\begin{enumerate}
\item \emph{Exact two-level factorization.}
Compute a coordinate Givens/QR factorization
\[
U=\Big(\prod_{k=1}^{K}\phi_{p_k,q_k}(V_k)\Big)\,D,
\qquad 1\le p_k<q_k\le N,\ \ V_k\in U(2),
\]
with $K\le N(N-1)/2$.
If desired, write $V_k=e^{i\theta_k}\widetilde V_k$ with $\widetilde V_k\in SU(2)$ and absorb the phases $e^{i\theta_k}$ into $D$.

\item \emph{Local approximation in $SU(2)$.}
Set $\delta:=\varepsilon/K$.
For each $k=1,\dots,K$, apply a Solovay--Kitaev routine to obtain a word $w_k$ over $G_2\cup G_2^{-1}$ such that
\[
\|\,\widetilde V_k-w_k\,\|\le \delta.
\]

\item \emph{Lift and concatenate.}
Define
\[
W:=\prod_{k=1}^{K}\phi_{p_k,q_k}(w_k),
\]
where $\phi_{p_k,q_k}(w_k)$ denotes the word obtained by replacing each letter of $w_k$ by its embedded generator in
$\phi_{p_k,q_k}(G_2)\cup \phi_{p_k,q_k}(G_2)^{-1}$.
\end{enumerate}
\end{quote}

\noindent
By Lemmas~\ref{lem:two_level_isometry} and~\ref{lem:error_accumulation}, the resulting pair $(W,D)$ satisfies $\|U-WD\|\le\varepsilon$.
\medskip

\begin{remark}[Interpretation in the geometric framework]
The factorization step may be viewed as selecting a finite sequence of logical two-planes (here, the coordinate planes
$W_{p_k,q_k}$) and prescribing on each plane a continuous $U(2)$ motion $V_k$. The Solovay--Kitaev step then replaces each
continuous $SU(2)$ motion by a short word over a fixed finite alphabet, and the lifting step transports these words to $U(N)$ through
the corresponding two-level embeddings. In this sense, the procedure separates a \emph{continuous} choice of support (a path through
logical qubits) from a \emph{discrete} approximation of the induced $SU(2)$ dynamics.
\end{remark}

\begin{remark}[Why the present framework is not merely ``QR + Solovay--Kitaev'']
\label{rem:not_just_QR_SK}
The QR/Givens factorization and the Solovay--Kitaev theorem are used here only as modular compilation primitives.
The conceptual contribution is orthogonal to their existence: we introduce an intrinsic \emph{descriptor layer} for elementary
operations in $U(2^n)$ based on the choice of a faithful embedded copy of $U(2)$ (equivalently, a point in the embedding landscape
$\Emb(U(2),U(N))$), together with the attendant gauge structure and geometric parameter spaces (notably the Grassmannian in the
two-level sector). The compilation pipeline then serves as an existence proof that finite alphabets can be lifted through these
intrinsic descriptors, rather than as the primary novelty.
\end{remark}

\subsection{Relation to fault-tolerant and hardware-constrained gate sets}
\label{subsec:ftqc_relation}

The finite-alphabet interface recorded in this section is deliberately \emph{agnostic} to the physical origin of the discrete gate set.
In fault-tolerant settings, the available alphabet is typically restricted (e.g.\ Clifford+$T$-type libraries, lattice-surgery primitives,
or other architecture-dependent native gates). Our framework separates (i) the \emph{intrinsic two-level lifting layer} in $U(N)$ from
(ii) the \emph{local} approximation engine in $SU(2)$: any routine that approximates target elements of $SU(2)$ over a prescribed finite
alphabet can be inserted into the pipeline, while the global error control after lifting through two-level embeddings remains unchanged.

Importantly, we do not claim optimal resource counts for any specific fault-tolerant metric (such as $T$-count). Rather, the point is that
hardware- or FTQC-imposed alphabets can be incorporated without altering the intrinsic descriptor layer or the two-level universality backbone,
thus providing a clean interface between $U(N)$-factorizations and architecture-specific synthesis.

%============================================================
\section{Structural backbone: the embedding landscape $\Emb(G,U(N))$ as a union of homogeneous manifolds}
\label{sec:EmbSU2UN}

\noindent
An embedding $\phi:G\hookrightarrow U(N)$ equips $\C^N$ with a unitary $G$-module structure, and the conjugation action of $U(N)$ on
$\Emb(G,U(N))$ is exactly unitary change of basis. Hence the geometry of $\Emb(G,U(N))$ is governed by the standard dichotomy:
\emph{representation type is discrete} (isotypic multiplicities), while \emph{basis choice is continuous} (unitary orbits).
We develop the structure at the $SU(2)$ level, where only finitely many types occur at fixed $N$; the case $G=U(2)$ adds determinant
twists and yields a typically countable refinement.

%------------------------------------------------------------
\subsection{Orbits, stabilizers, and multiplicity data for $SU(2)$}
\label{subsec:Emb_SU2_orbits_concise}
%------------------------------------------------------------

Let $\{V_d\}_{d\in\N_0}$ denote the irreducible unitary $SU(2)$-modules, where $\dim_\C(V_d)=d+1$.
Since $SU(2)$ is compact, every finite-dimensional unitary representation is completely reducible and is uniquely determined up to
unitary equivalence by its irreducible multiplicities; see, e.g., \cite{BrockerDieck1985,Hall2015,Knapp2002}.
Thus any $\phi\in\Emb(SU(2),U(N))$ yields a finite-support multiplicity family $m=(m_d)_{d\ge0}$ such that
\begin{equation}\label{eq:isotypic_decomposition_SU2_concise}
\C^N \ \cong\ \bigoplus_{d\ge0}\big(\C^{m_d}\otimes V_d\big),
\qquad
\sum_{d\ge0} m_d(d+1)=N.
\end{equation}

\begin{lemma}[Orbits are homogeneous]\label{lem:orbit_homogeneous_SU2_concise2}
For $\phi\in\Emb(SU(2),U(N))$, the $U(N)$-orbit
$\mathcal O_\phi:=\{W\phi(\cdot)W^{-1}:W\in U(N)\}$ is a smooth homogeneous manifold and
\[
\mathcal O_\phi \ \cong\ U(N)\big/ Z_{U(N)}\!\big(\phi(SU(2))\big),
\]
where $Z_{U(N)}(\phi(SU(2)))=\{W\in U(N):W\phi(g)=\phi(g)W\ \forall g\in SU(2)\}$.
\end{lemma}

\begin{proof}
Consider the smooth action of $U(N)$ on $\Emb(SU(2),U(N))$ by conjugation,
$(W\cdot \phi)(g)=W\phi(g)W^{-1}$.
The stabilizer of $\phi$ consists precisely of those $W$ commuting with $\phi(SU(2))$, i.e.\ $Z_{U(N)}(\phi(SU(2)))$.
By the orbit--stabilizer theorem for Lie group actions, the orbit is a smooth immersed submanifold and is diffeomorphic to the
homogeneous space $U(N)/Z_{U(N)}(\phi(SU(2)))$; see, e.g., \cite{Hall2015,Knapp2002}.
\end{proof}

\begin{lemma}[Centralizer and orbit dimension]\label{lem:centralizer_SU2_concise2}
Let $\phi\in\Emb(SU(2),U(N))$ have multiplicities $m=(m_d)_{d\ge0}$ as in \eqref{eq:isotypic_decomposition_SU2_concise}.
Then
\[
Z_{U(N)}\!\big(\phi(SU(2))\big)\ \cong\ \prod_{d\ge0} U(m_d),
\qquad
\dim(\mathcal O_\phi)=N^2-\sum_{d\ge0} m_d^2.
\]
\end{lemma}

\begin{proof}
Write $\C^N\simeq\bigoplus_{d\ge0}(\C^{m_d}\otimes V_d)$ so that $\phi(g)=\bigoplus_{d}(I_{m_d}\otimes \rho_d(g))$.
If $T\in\End(\C^N)$ commutes with all $\phi(g)$, then $T$ preserves each isotypic component and, on $\C^{m_d}\otimes V_d$,
Schur's lemma gives $T|_{\C^{m_d}\otimes V_d}=A_d\otimes I_{V_d}$ for some $A_d\in M_{m_d}(\C)$; hence
$\End_{SU(2)}(\C^N)\cong\bigoplus_{d\ge0}(M_{m_d}(\C)\otimes I_{V_d})$.
Intersecting with $U(N)$ forces each $A_d$ to be unitary, yielding
$Z_{U(N)}(\phi(SU(2)))\cong\prod_{d\ge0}U(m_d)$.
Finally, $\dim_\R U(N)=N^2$ and $\dim_\R U(m_d)=m_d^2$, so
$\dim(\mathcal O_\phi)=N^2-\sum_{d\ge0}m_d^2$.
\end{proof}

%------------------------------------------------------------
\subsection{Faithfulness for $SU(2)$ embeddings}
\label{subsec:Emb_SU2_faithfulness_concise}
%------------------------------------------------------------

\begin{lemma}[Parity criterion]\label{lem:faithfulness_SU2_parity_concise2}
Let $\rho=\bigoplus_{d\ge0}(\C^{m_d}\otimes V_d)$ be a unitary representation of $SU(2)$.
Then $\rho$ is faithful if and only if $m_d>0$ for some odd $d$.
\end{lemma}

\begin{proof}
The only nontrivial proper normal subgroup of $SU(2)$ is its center $\{\pm I\}$.
On $V_d$, the element $-I$ acts by $(-1)^d$, hence $\rho(-I)=I$ iff all occurring $d$ are even. Therefore $\ker(\rho)=\{I\}$ iff some
odd $d$ occurs.
\end{proof}

%------------------------------------------------------------
\subsection{Homogeneous stratification of $\Emb(SU(2),U(N))$}
\label{subsec:Emb_SU2_stratification_concise}
%------------------------------------------------------------

\begin{theorem}[Homogeneous decomposition of $\Emb(SU(2),U(N))$]\label{thm:Emb_SU2_decomposition_main_concise2}
Fix $N\in\N$. Let $\mathscr F_{SU(2)}(N)$ be the set of finite-support families $m=(m_d)_{d\ge0}$ with $m_d\in\N_0$ such that
\[
\sum_{d\ge0} m_d(d+1)=N
\quad\text{and}\quad
\exists\, d \text{ odd with } m_d>0.
\]
Then
\[
\Emb(SU(2),U(N))=\bigsqcup_{m\in\mathscr F_{SU(2)}(N)} \mathcal O_m,
\qquad
\mathcal O_m \cong U(N)\Big/\Big(\prod_{d\ge0}U(m_d)\Big),
\]
and $\dim(\mathcal O_m)=N^2-\sum_{d\ge0}m_d^2$.
\end{theorem}

\begin{proof}
By complete reducibility, every unitary homomorphism $SU(2)\to U(N)$ decomposes into irreducibles as in
\eqref{eq:isotypic_decomposition_SU2_concise}; the multiplicities $m=(m_d)_{d\ge0}$ determine the representation up to unitary
equivalence \cite{BrockerDieck1985,Hall2015,Knapp2002}. In particular, two homomorphisms are $U(N)$-conjugate if and only if they have
the same multiplicities, so the $U(N)$-orbits are indexed by families $m$ satisfying the dimension constraint.

The embedding condition is faithfulness, which for $SU(2)$ is equivalent to the parity requirement in
Lemma~\ref{lem:faithfulness_SU2_parity_concise2}. For each admissible $m$, choose an embedding $\phi_m$ of type $m$ and set
$\mathcal O_m:=U(N)\cdot \phi_m$. Lemma~\ref{lem:orbit_homogeneous_SU2_concise2} gives
$\mathcal O_m\cong U(N)/Z_{U(N)}(\phi_m(SU(2)))$, and Lemma~\ref{lem:centralizer_SU2_concise2} identifies the stabilizer with
$\prod_{d\ge0}U(m_d)$ and yields $\dim(\mathcal O_m)=N^2-\sum_{d\ge0}m_d^2$.
\end{proof}

\begin{corollary}[Finiteness of strata]\label{cor:finiteness_strata_SU2_concise2}
For fixed $N$, the set $\mathscr F_{SU(2)}(N)$ is finite. Hence $\Emb(SU(2),U(N))$ is a finite disjoint union of $U(N)$-orbits.
\end{corollary}

\begin{remark}[The two-level stratum]\label{rem:two_level_stratum_concise}
The two-level sector corresponds to $m_1=1$, $m_0=N-2$, $m_d=0$ otherwise, i.e.\ $V_1\oplus V_0^{\oplus (N-2)}$.
It is the orbit with stabilizer $U(1)\times U(N-2)$ and dimension $N^2-(1+(N-2)^2)=4N-5$, and it is organized by the
Grassmannian model developed in Section~\ref{sec:logical_qubits}.
\end{remark}

%------------------------------------------------------------
\subsection{Remark on $U(2)$ and determinant twists}
\label{subsec:Emb_U2_det_twists_concise}
%------------------------------------------------------------

\begin{remark}[$U(2)$ adds determinant twists]\label{rem:SU2_vs_U2_det_twists_concise2}
Every irreducible unitary $U(2)$-representation is of the form $\Sym^{d}(\C^2)\otimes \det^{k}$ with $d\in\N_0$ and $k\in\Z$, and has
dimension $d+1$ \cite{BrockerDieck1985,Hall2015,Knapp2002}. Hence a unitary $U(2)$-representation decomposes as
$\bigoplus_{(d,k)}(\C^{m_{d,k}}\otimes \Sym^{d}(\C^2)\otimes \det^{k})$ with $\sum_{(d,k)}m_{d,k}(d+1)=N$.
Restricting to $SU(2)\subset U(2)$ forgets the twist $k$ and retains only $m_d=\sum_k m_{d,k}$, recovering the $SU(2)$ multiplicity data.
\end{remark}

\begin{example}[Small-$N$ illustrations]\label{ex:smallN_strata_concise}
The stratification in Theorem~\ref{thm:Emb_SU2_decomposition_main_concise2} is already visible for $N=4$ and $N=8$.

\smallskip
\noindent\textbf{(i) $N=4$.}
There are (at least) two distinct strata in $\Emb(SU(2),U(4))$:
\begin{enumerate}
\item \emph{Two-level type} $m_1=1$, $m_0=2$ (i.e.\ $V_1\oplus V_0^{\oplus 2}$). Here the stabilizer is
$U(1)\times U(2)$ and $\dim(\mathcal O_m)=16-(1^2+2^2)=11$. This is the Grassmannian sector modeled in
Section~\ref{sec:logical_qubits}.
\item \emph{Irreducible type} $m_3=1$ (i.e.\ $V_3\simeq \Sym^3(\C^2)$, spin-$\tfrac32$). Here the stabilizer is $U(1)$ and
$\dim(\mathcal O_m)=16-1=15$. In particular, the action has no nontrivial invariant subspace, so it does not single out a
distinguished two-plane support.
\end{enumerate}

\smallskip
\noindent\textbf{(ii) $N=8$.}
Two extremal (and computationally interpretable) strata are:
\begin{enumerate}
\item \emph{Two-level type} $m_1=1$, $m_0=6$ (i.e.\ $V_1\oplus V_0^{\oplus 6}$). Stabilizer $U(1)\times U(6)$ and
$\dim(\mathcal O_m)=64-(1^2+6^2)=27$.
\item \emph{Irreducible type} $m_7=1$ (i.e.\ $V_7$, dimension $8$ irreducible). Stabilizer $U(1)$ and
$\dim(\mathcal O_m)=64-1=63$.
\end{enumerate}

\smallskip
\noindent
These cases illustrate the general mechanism: the representation type (multiplicity data) determines the stabilizer
$\prod_d U(m_d)$ and hence the orbit dimension $N^2-\sum_d m_d^2$. The two-level strata are precisely the ones that interface
directly with the Grassmannian logical-qubit model, whereas irreducible strata correspond to intrinsically higher-dimensional logical
primitives.
\end{example}

%============================================================
\section{Supporting layer: a variational (metric) cost for elementary gates}
\label{sec:variational_principle}

\noindent
We equip $U(N)$ with the bi-invariant Riemannian metric induced by the Hilbert--Schmidt pairing on $\mathfrak u(N)$.
For any embedding $\phi\in\Emb(G,U(N))$ (with $G\in\{SU(2),U(2)\}$), the image $K=\phi(G)$ is a closed Lie subgroup of $U(N)$.
Bi-invariance implies that $K$ is totally geodesic, so constrained minimizers inside $K$ coincide with ambient geodesics whose initial
velocity lies in the embedded Lie algebra $\mathfrak{k}=d\phi_e(\Lie(G))$.
As a consequence, minimal-energy implementations in $K$ are precisely constant-speed one-parameter subgroups generated by
minimal-norm logarithms in $\mathfrak{k}$.

\begin{remark}[On $U(2)$ embeddings and phase degrees of freedom]\label{rem:U2_vs_SU2_variational}
If $\widehat\phi\in\Emb(U(2),U(N))$ and $\phi:=\widehat\phi|_{SU(2)}$, then $\widehat\phi(U(2))$ contains $\phi(SU(2))$ and an
additional commuting $U(1)$ subgroup coming from the center of $U(2)$.
At the Lie-algebra level $\mathfrak u(2)=\mathfrak{su}(2)\oplus i\R I$ is a direct sum of commuting ideals; the pullback of the
Hilbert--Schmidt product via $d\widehat\phi_e$ is $\Ad$-invariant and has no cross-term between these ideals.
Accordingly, the non-abelian $SU(2)$ and central $U(1)$ contributions to constrained geodesics and minimal energy can be treated
separately when needed.
\end{remark}


%------------------------------------------------------------
\subsection{Hilbert--Schmidt metric and basic variational functionals}
\label{subsec:HS_metric_functionals}
%------------------------------------------------------------

The Hilbert--Schmidt pairing on $\mathfrak u(N)$,
\[
\langle X,Y\rangle_{\mathrm{hs}}:=\frac12\Tr(X^\dagger Y),
\]
is $\Ad$-invariant and induces a bi-invariant Riemannian metric on $U(N)$ by left/right translation.
For an absolutely continuous curve $U:[0,1]\to U(N)$ define the body velocity
\[
A(t):=\dot U(t)\,U(t)^{-1}\in\mathfrak{u}(N)\quad\text{for a.e.\ }t\in[0,1],
\]
and the energy and length
\begin{align}
\mathcal{E}[U]
&:=\frac12\int_0^1 \|A(t)\|_{\mathrm{hs}}^2\,dt,
\label{eq:energy_functional}
\\
\mathcal{L}[U]
&:=\int_0^1 \|A(t)\|_{\mathrm{hs}}\,dt.
\label{eq:length_functional}
\end{align}
As usual, among curves with fixed endpoints, minimizing $\mathcal{E}$ is equivalent to minimizing $\mathcal{L}$ after
constant-speed reparametrization.

%------------------------------------------------------------
\subsection{Bi-invariant geodesics and totally geodesic subgroups}
\label{subsec:geodesics_and_totally_geodesic}
%------------------------------------------------------------

\begin{lemma}[Geodesics are one-parameter subgroups]\label{lem:geodesics_one_parameter}
Let $G$ be a Lie group equipped with a bi-invariant Riemannian metric induced by an $\Ad$-invariant inner product on $\mathfrak g$.
Then the geodesics in $G$ are exactly the curves
\[
\gamma(t)=g_0\exp(tX),\qquad g_0\in G,\ X\in\mathfrak g.
\]
\end{lemma}

\begin{proof}
For a bi-invariant metric, left and right translations are isometries.
The Levi--Civita connection satisfies $\nabla_{X^R}Y^R=\tfrac12[X,Y]^R$ for right-invariant fields.
Hence the geodesic equation is equivalent to constant body velocity.
Thus $\gamma(t)=g_0\exp(tX)$. See \cite[Ch.~4]{Hall2015} or \cite[\S~II.1]{Knapp2002}.
\end{proof}

\begin{lemma}[Lie subgroups are totally geodesic]\label{lem:embedded_subgroup_totally_geodesic}
Let $G$ be a Lie group with a bi-invariant Riemannian metric and let $H\subset G$ be a Lie subgroup with the induced metric.
Then $H$ is totally geodesic in $G$.
Equivalently, if $\gamma$ is a $G$-geodesic with $\gamma(0)\in H$ and $\dot\gamma(0)\in T_{\gamma(0)}H$, then $\gamma(t)\in H$ for all $t$.
\end{lemma}

\begin{proof}
By Lemma~\ref{lem:geodesics_one_parameter}, $\gamma(t)=\gamma(0)\exp(tX)$ for some $X\in\mathfrak g$.
If $\gamma(0)\in H$ and $\dot\gamma(0)\in T_{\gamma(0)}H$, then the right-translated velocity $X=\dot\gamma(0)\gamma(0)^{-1}$
lies in $\mathfrak h=\Lie(H)$. Hence $\exp(tX)\in H$ and therefore $\gamma(t)\in H$ for all $t$.
\end{proof}

\begin{remark}[Application to embedded copies of $G\in\{SU(2),U(2)\}$]
\label{rem:totally_geodesic_embedded_SU2}
Fix $G\in\{SU(2),U(2)\}$ and $\phi\in\Emb(G,U(N))$. Then $K=\phi(G)$ is a closed Lie subgroup of $U(N)$ with Lie algebra
\[
\mathfrak{k}=d\phi_e(\Lie(G)).
\]
Since the Hilbert--Schmidt metric on $U(N)$ is bi-invariant, Lemma~\ref{lem:embedded_subgroup_totally_geodesic} implies that $K$ is
totally geodesic.
\end{remark}

%------------------------------------------------------------
\subsection{A constrained variational principle on an embedded subgroup}
\label{subsec:constrained_variational_principle}
%------------------------------------------------------------

\begin{definition}[Admissible curves constrained to an embedded subgroup]\label{def:admissible_curves}
Fix $G\in\{SU(2),U(2)\}$, $\phi\in\Emb(G,U(N))$, and $K:=\phi(G)$. For $U_{\star}\in K$ define
\[
\mathcal{A}(I,U_{\star};K)
:=\big\{\,U\in AC([0,1],K):\ U(0)=I,\ U(1)=U_{\star}\,\big\}.
\]
\end{definition}

\paragraph{Minimal logarithms.}
For $U_{\star}\in K$ set
\[
\Log_{\mathfrak{k}}(U_{\star}):=\{\,X\in\mathfrak{k}:\ \exp(X)=U_{\star}\,\}.
\]
Since $K$ is compact and connected, $\Log_{\mathfrak{k}}(U_{\star})\neq\varnothing$.

\begin{theorem}[Constrained energy minimizers are one-parameter subgroups]
\label{thm:variational_characterization}
Fix $G\in\{SU(2),U(2)\}$ and $\phi\in\Emb(G,U(N))$, and set $K=\phi(G)$ with Lie algebra $\mathfrak{k}=d\phi_e(\Lie(G))$.
Equip $U(N)$ with the Hilbert--Schmidt bi-invariant metric and $K$ with the induced metric. Fix $U_{\star}\in K$. Then:
\begin{enumerate}
\item[(i)] There exists at least one minimizer of $\mathcal{E}$ over $\mathcal{A}(I,U_{\star};K)$.
\item[(ii)] A curve $U\in\mathcal{A}(I,U_{\star};K)$ minimizes $\mathcal{E}$ if and only if it is a constant-speed minimizing geodesic in $K$.
Equivalently, every minimizer is of the form
\begin{equation}\label{eq:minimizer_form}
U(t)=\exp(tX),\qquad
X\in\Log_{\mathfrak{k}}(U_{\star}),
\qquad
\|X\|_{\mathrm{hs}}=\min_{Y\in\Log_{\mathfrak{k}}(U_{\star})}\|Y\|_{\mathrm{hs}}.
\end{equation}
\item[(iii)] Every minimizer \eqref{eq:minimizer_form} is also an ambient geodesic of $U(N)$; conversely, any ambient geodesic from $I$
to $U_{\star}$ with initial velocity in $\mathfrak{k}$ is a constrained minimizer.
\item[(iv)] The minimal energy value is
\begin{equation}\label{eq:Emin_log}
\min_{U\in\mathcal{A}(I,U_{\star};K)}\mathcal{E}[U]
=\frac12\min_{X\in\Log_{\mathfrak{k}}(U_{\star})}\|X\|_{\mathrm{hs}}^{\,2}.
\end{equation}
\end{enumerate}
\end{theorem}

\begin{proof}
(i) The subgroup $K$ is compact as the continuous image of the compact group $G$ and is a smooth Riemannian manifold with the induced metric.
By Hopf--Rinow, any two points in a compact Riemannian manifold are joined by a minimizing geodesic; hence an energy minimizer exists
\cite{doCarmoRiemannian,KobayashiNomizuI}.

(ii) On any Riemannian manifold, energy minimizers with fixed endpoints are exactly constant-speed minimizing geodesics
\cite[\S~6.4]{doCarmoRiemannian}. Since the metric on $K$ is the restriction of a bi-invariant metric, it is bi-invariant on $K$.
Thus, by Lemma~\ref{lem:geodesics_one_parameter}, geodesics in $K$ are one-parameter subgroups $t\mapsto \exp(tX)$ with $X\in\mathfrak{k}$.
The endpoint condition forces $\exp(X)=U_\star$, i.e.\ $X\in\Log_{\mathfrak{k}}(U_{\star})$.
For such a curve, the body velocity is constant and equals $X$, hence $\mathcal{E}[\exp(tX)]=\tfrac12\|X\|_{\mathrm{hs}}^{2}$.
Minimizing $\mathcal{E}$ is therefore equivalent to choosing a minimal-norm logarithm in $\Log_{\mathfrak{k}}(U_\star)$.

(iii) By Lemma~\ref{lem:embedded_subgroup_totally_geodesic}, $K$ is totally geodesic in $U(N)$, so every $K$-geodesic is an ambient
$U(N)$-geodesic. Conversely, any ambient geodesic from $I$ with initial velocity in $\mathfrak{k}$ remains in $K$ and is therefore a
constrained geodesic, hence an energy minimizer by (ii).

(iv) This follows by evaluating $\mathcal{E}$ along $t\mapsto \exp(tX)$ and taking the minimum over $X\in\Log_{\mathfrak{k}}(U_{\star})$,
using (ii).
\end{proof}

\begin{remark}[Pullback metric and representation-dependent scaling]\label{rem:metric_scaling_index}
The induced inner product on $\mathfrak{k}$ pulls back via $d\phi_e$ to an $\Ad$-invariant inner product on $\mathfrak g=\Lie(G)$.
If $G=SU(2)$, then $\mathfrak g=\mathfrak{su}(2)$ is simple, so any $\Ad$-invariant inner product is a positive scalar multiple of a fixed
normalization; equivalently, there exists $c_\phi>0$ such that
\[
\big\langle d\phi_e(A),d\phi_e(B)\big\rangle_{\mathrm{hs}}=c_\phi\,\langle A,B\rangle_{\mathrm{ref}}
\qquad (A,B\in\mathfrak{su}(2)).
\]
If $G=U(2)$, then $\mathfrak{u}(2)=\mathfrak{su}(2)\oplus i\R I$ is reductive and $\Ad$-invariant inner products are determined by two
positive parameters, one on each ideal.
In the two-level sector of Section~\ref{sec:logical_qubits}, the natural normalization inherited from $\mathfrak u(N)$ yields $c_\phi=1$.
\end{remark}

%------------------------------------------------------------
\subsection{Explicit evaluation in $U(2)$ and $SU(2)$}
\label{subsec:explicit_U2_SU2}
%------------------------------------------------------------

Since $SU(2)$ (and $U(2)$) logarithms admit explicit eigen-angle formulae, we record a convenient reference statement.

\begin{lemma}[Minimal Hilbert--Schmidt logarithms in $U(2)$ and $SU(2)$]\label{lem:minlog_U2_eigenangles}
Let $V\in U(2)$. Choose a unitary diagonalization
\[
V = Q\,\diag\!\big(e^{i\theta_1},e^{i\theta_2}\big)\,Q^{-1},
\qquad \theta_1,\theta_2\in(-\pi,\pi],
\]
(where $\theta_j$ are principal eigen-angles, uniquely determined up to permutation).
Then
\[
\Log_{\mathfrak{u}(2)}(V)
=
\left\{
Q\,\diag\!\big(i(\theta_1+2\pi k_1),\,i(\theta_2+2\pi k_2)\big)\,Q^{-1}
:\ k_1,k_2\in\mathbb{Z}
\right\},
\]
and
\begin{equation}\label{eq:minlog_U2_value}
\min_{X\in\Log_{\mathfrak{u}(2)}(V)}\|X\|_{\mathrm{hs}}^2
\;=\;
\frac12\big(\theta_1^2+\theta_2^2\big).
\end{equation}
In particular, if $V\in SU(2)$ has eigenvalues $e^{\pm i\alpha}$ with $\alpha\in[0,\pi]$, then
\[
\min_{X\in\Log_{\mathfrak{su}(2)}(V)}\|X\|_{\mathrm{hs}}=\alpha,
\qquad
\min_{X\in\Log_{\mathfrak{su}(2)}(V)}\|X\|_{\mathrm{hs}}^{\,2}=\alpha^2.
\]
\end{lemma}

\begin{proof}
The Hilbert--Schmidt norm is unitarily invariant, hence we may conjugate by $Q^{-1}$ and reduce to the diagonal case.
Diagonal logarithms of $\diag(e^{i\theta_1},e^{i\theta_2})$ are exactly $\diag(i(\theta_1+2\pi k_1),\,i(\theta_2+2\pi k_2))$.
Their squared Hilbert--Schmidt norm equals
$\tfrac12\big((\theta_1+2\pi k_1)^2+(\theta_2+2\pi k_2)^2\big)$, minimized by choosing the representatives closest to $0$, i.e.\ by
$k_1=k_2=0$ for principal angles. This yields \eqref{eq:minlog_U2_value}. The $SU(2)$ case follows from
$(\theta_1,\theta_2)=(\alpha,-\alpha)$.
\end{proof}

\begin{remark}[Non-uniqueness and cut locus]\label{rem:nonuniqueness_log}
The minimizer in Theorem~\ref{thm:variational_characterization} need not be unique: non-uniqueness occurs precisely when
$\Log_{\mathfrak{k}}(U_{\star})$ contains more than one element of minimal Hilbert--Schmidt norm, equivalently when $U_{\star}$ lies in
the cut locus of $I$ in $K$. For generic $U_{\star}$ the minimal logarithm is unique.
\end{remark}

%------------------------------------------------------------
\subsection{Interpretation and the two-level sector}
\label{subsec:interpretation_minimal_action}
%------------------------------------------------------------

Theorem~\ref{thm:variational_characterization} identifies the minimal-action implementation of a target $U_{\star}\in K=\phi(G)$:
among all admissible curves constrained to remain in $K$, the minimizers are constant-speed geodesics
$t\mapsto\exp(tX)$ generated by minimal-norm logarithms $X\in\mathfrak{k}$.

In the two-level (Grassmannian) sector of Section~\ref{sec:logical_qubits}, generators take the descriptor form
$X=X_{W,f}(A)=fAf^{-1}\oplus 0$ with $A\in\mathfrak{su}(2)$, and Lemma~\ref{lem:HS_norm_two_level_generator} implies that
$\|X\|_{\mathrm{hs}}=\|A\|_{\mathrm{hs}}$ is frame-independent.
Thus the minimal-action cost of a two-level elementary gate is exactly the minimal Hilbert--Schmidt norm of its internal $2\times 2$
generator, providing a direct bridge between intrinsic geometry on $U(N)$ and concrete $SU(2)$ control parameters.
When an additional abelian $U(1)$ phase is tracked (i.e.\ $SU(N)[W]$ extended to $U(N)[W]$), its contribution is generated by a commuting
central direction and can be compiled/penalized separately, consistent with Remark~\ref{rem:U2_vs_SU2_variational}.

\section{Conclusions and final remarks}
\label{sec:conclusion}
%============================================================

We introduced an intrinsic notion of elementary gates in $U(2^n)$ based on Lie group embeddings: a two-level logical degree of freedom
is modeled by an internal copy of $SU(2)$ inside $U(N)$, $N=2^n$ (and, when phases are kept explicit, by an embedding of $U(2)$).
This viewpoint separates \emph{discrete} representation data from \emph{continuous} gauge/basis choices, and yields a structured,
basis-invariant parameter space of admissible primitives.

\smallskip
\noindent\emph{Embedding landscape and a canonical sector.}
At fixed $N$, the space $\Emb(SU(2),U(N))$ decomposes into finitely many $U(N)$-conjugacy orbits, each an explicit homogeneous manifold
indexed by isotypic multiplicities, with stabilizers computed as centralizers.
Among these strata, the two-level sector provides the canonical computational interface: a logical qubit is a two-plane
$W\in\Gr_2(\C^N)$, and the corresponding subgroup is $SU(N)[W]$.
Given any Stiefel frame $F\in\C^{N\times 2}$ with $F^\dagger F=I_2$ and $\Ran(F)=W$, two-level gates admit the concrete realization
\begin{equation}\label{eq:conclusion_two_level_realization}
U_{W,F}(S)=I_N+F(S-I_2)F^\dagger\in SU(N)[W],\qquad S\in SU(2),
\end{equation}
and similarly with $S$ replaced by $V\in U(2)$ for $U(N)[W]$.
The support subgroup depends only on $W$, while the internal labeling depends on $F$ only up to the natural $PSU(2)$ gauge.

\smallskip
\noindent\emph{Expressiveness and a finite-alphabet interface.}
The universality mechanism developed in Section~\ref{sec:universality} shows that two-level $SU(2)$ gates already generate $SU(N)$, and
that full reachability in $U(N)$ follows after adjoining global phase (equivalently, by explicit bookkeeping of abelian factors).
Section~\ref{sec:SK} then provides a modular compilation route: exact two-level factorizations reduce synthesis in $U(N)$ to repeated
approximation problems in $SU(2)$ (after optional phase normalization), and finite-alphabet approximation on $SU(2)$ lifts through
two-level embeddings to uniform operator-norm control in $U(N)$, with accuracy--length guarantees inherited from the underlying
$SU(2)$ approximation scheme.

\smallskip
\noindent\emph{Limitations and directions.}
Several mathematically natural extensions emerge:
\begin{itemize}
\item \emph{Beyond bi-invariant costs:} the Hilbert--Schmidt metric yields clean variational statements (total geodesy and
minimal-logarithm characterizations), whereas realistic control costs are typically anisotropic; extending the analysis to
right-invariant or sub-Riemannian models is a natural next step.
\item \emph{Descriptor selection:} a fixed $U\in U(N)$ may admit multiple embedding-elementary descriptions across different strata and
supports; choosing descriptors optimally leads to constrained optimization on stratified spaces.
\item \emph{Depth/complexity refinement:} QR/Givens-type synthesis is explicit but not depth-optimal in general; sharper complexity bounds
should exploit target structure (symmetry, sparsity, tensor-network structure) and may benefit from structure-adapted factorizations and
improved $SU(2)$ approximation schemes.
\end{itemize}


%\paragraph{Acknowledgments}
%We are grateful for the technical assistance of A. Author.

\paragraph{Funding statement}
This work was supported by the Generalitat Valenciana under grant COMCUANTICA/007 (QUANTWin), by the Agreement between the Directorate-General for Innovation of the Ministry of Innovation, Industry, Trade and Tourism of the Generalitat Valenciana and the Universidad CEU Cardenal Herrera, and by Universidad CEU Cardenal Herrera under grants INDI25/17 and GIR25/14.

\paragraph{Competing interests}
None

\paragraph{Data availability statement}
No datasets were generated or analyzed during the current study.

\paragraph{Ethical standards}
The research meets all ethical guidelines, including adherence to the legal requirements of the study country.

\paragraph{Author contributions}
All authors contributed equally to this work. All authors jointly designed the study, developed the methodology, and drafted the manuscript. All authors revised the manuscript critically and approved the final version.

%\paragraph{Supplementary material}
%State whether any supplementary material intended for publication has been provided with the submission.


\begin{thebibliography}{99}

\bibitem{BarencoEtAl1995}
A.~Barenco, C.~H.~Bennett, R.~Cleve, D.~P.~DiVincenzo, N.~Margolus, P.~Shor, T.~Sleator, J.~A.~Smolin, and H.~Weinfurter,
\newblock ``Elementary gates for quantum computation,''
\newblock \emph{Physical Review A} \textbf{52} (1995), 3457--3467.

\bibitem{BrockerDieck1985}
T.~Br{\"o}cker and T.~tom Dieck,
\newblock \emph{Representations of Compact Lie Groups},
\newblock Graduate Texts in Mathematics, Vol.~98, Springer, Berlin--Heidelberg, 1985.

\bibitem{DawsonNielsen2006}
C.~M.~Dawson and M.~A.~Nielsen,
\newblock The Solovay--Kitaev algorithm,
\newblock \emph{Quantum Information \& Computation} \textbf{6}(1) (2006), 81--95.
\newblock (Preprint: arXiv:quant-ph/0505030.)

\bibitem{DeutschBarencoEkert1995}
D.~Deutsch, A.~Barenco, and A.~Ekert,
\newblock ``Universality in quantum computation,''
\newblock \emph{Proceedings of the Royal Society of London A} \textbf{449} (1995), 669--677.

\bibitem{DiVincenzo1995}
D.~P.~DiVincenzo,
\newblock ``Two-bit gates are universal for quantum computation,''
\newblock \emph{Physical Review A} \textbf{51} (1995), 1015--1022.

\bibitem{doCarmoRiemannian}
M.~P.~do~Carmo,
\newblock \emph{Riemannian Geometry},
\newblock Mathematics: Theory \& Applications,
Birkh{\"a}user, Boston, 1992.

\bibitem{DowlingNielsen2008}
M.~R.~Dowling and M.~A.~Nielsen,
\newblock ``The geometry of quantum computation,''
\newblock \emph{Quantum Information \& Computation} \textbf{8} (2008), no.~10, 861--899.

\bibitem{GolubVanLoan2013}
G.~H.~Golub and C.~F.~Van Loan,
\newblock \emph{Matrix Computations},
\newblock 4th ed., Johns Hopkins Studies in the Mathematical Sciences,
Johns Hopkins University Press, Baltimore, 2013.

\bibitem{Hall2015}
B.~C.~Hall,
\newblock \emph{Lie Groups, Lie Algebras, and Representations: An Elementary Introduction},
\newblock 2nd ed., Graduate Texts in Mathematics, Vol.~222, Springer, Cham, 2015.

\bibitem{HarrowRechtChuang2002}
A.~W.~Harrow, B.~Recht, and I.~L.~Chuang,
\newblock ``Efficient discrete approximations of quantum gates,''
\newblock \emph{Journal of Mathematical Physics} \textbf{43} (2002), no.~9, 4445--4451.

\bibitem{Higham2008}
N.~J.~Higham,
\newblock \emph{Functions of Matrices: Theory and Computation},
\newblock Society for Industrial and Applied Mathematics (SIAM), Philadelphia, PA, 2008.
\newblock \textsc{doi}:~10.1137/1.9780898717778.

\bibitem{HornJohnson2013}
R.~A.~Horn and C.~R.~Johnson,
\newblock \emph{Matrix Analysis},
\newblock 2nd ed., Cambridge University Press, Cambridge, 2013.

\bibitem{Kitaev2002}
A.~Y.~Kitaev, A.~H.~Shen, and M.~N.~Vyalyi,
\newblock \emph{Classical and Quantum Computation},
\newblock Graduate Studies in Mathematics, Vol.~47,
American Mathematical Society, Providence, RI, 2002.

\bibitem{Knapp2002}
A.~W.~Knapp,
\newblock \emph{Lie Groups Beyond an Introduction},
\newblock 2nd ed., Progress in Mathematics, Vol.~140, Birkh{\"a}user, Boston, 2002.

\bibitem{KobayashiNomizuI}
S.~Kobayashi and K.~Nomizu,
\newblock \emph{Foundations of Differential Geometry, Vol.~I},
\newblock Wiley Classics Library,
% (Original edition: 1963.)

\bibitem{Kuperberg2025BreakingCubicSK}
G.~Kuperberg,
\newblock Breaking the cubic barrier in the Solovay--Kitaev algorithm,
\newblock \emph{arXiv preprint} arXiv:2306.13158v2 (Oct.~8, 2025).

\bibitem{NielsenChuang2010}
M.~A.~Nielsen and I.~L.~Chuang,
\newblock \emph{Quantum Computation and Quantum Information: 10th Anniversary Edition},
\newblock Cambridge University Press, Cambridge, 2010.

\bibitem{NielsenDowlingGuDoherty2006}
M.~A.~Nielsen, M.~R.~Dowling, M.~Gu, and A.~C.~Doherty,
\newblock ``Quantum computation as geometry,''
\newblock \emph{Science} \textbf{311} (2006), no.~5764, 1133--1135.

\bibitem{ReckZeilingerBernsteinBertani1994}
M.~Reck, A.~Zeilinger, H.~J.~Bernstein, and P.~Bertani,
\newblock ``Experimental realization of any discrete unitary operator,''
\newblock \emph{Physical Review Letters} \textbf{73} (1994), 58--61.

\end{thebibliography}


\end{document}
